\chapter{Anatomie und Physiologie}
\label{chp:ANA}

\emph{Maintainer: Sebastian Gabriel}

 
\gAuthNotes{
Changelog: 

\begin{description}
    \item[2009-02-25] Start Changelog.
    \item[2009-05-06] GABS: Korrekturen von KOCH
    eingearbeitet.
\end{description}
}

\minitoc

\section{Die Zelle}\index{Zelle}
\gLayout{0}{\gTxBeschreibung}
{
Die Zelle ist die kleinste lebende Einheit in unserem
Körper. Sie besteht u.a. aus einem
\gE{Zellkern}\index{Zellkern}, der die gesamten Erbinformationen in Form von DNA
enthält und verschiedenen Zellorganellen\index{Zellorganellen}. Sie wird von der
Zellmembran\index{Zellmembran} umhüllt.
}
{
\begin{itemize}
    \item Kleinste lebende Einheit.
    \item Zellkern
    \begin{itemize}
        \item Erbinformationen (DNA)
    \end{itemize}
    \item Zellorganellen
\end{itemize}
}
{}{}{}



\begin{figure}[H]
    \includegraphics[width=\linewidth]{./images/anatomie/Animal_cell_structure_de-edited.png}
    \gCaptionof{figure}{\gAuthNotesPicLegalOk{}{PD}Eine
    Zelle, Schema. \gSourcePicture{Mariana Ruiz Villarreal,
    Public domain}}
\end{figure}

\gMarriedNG{Gewebe}
{
Artverwandte Zellen können einen
Zellverband bilden, welcher \gT{Gewebe}\index{Gewebe} genannt wird. Diese
Gewebe können dann größere Strukturen, wie z.B. Organe, "`erbauen"'.
}
{
\begin{itemize}
    \item Zellverband.
    \item → größere Strukturen.
\end{itemize}
}


\section{Das Skelett und der Bewegungsapparat}
\index{Skelett}

\begin{figure}[H]
    \includegraphics[width=10cm]{./images/anatomie/Human_skeleton_front_de.png} \gCaptionof{figure}{\gAuthNotesPicLegalOk{}{ PD }Das Menschliche Skelett.
    \gSourcePicture{Mariana Ruiz Villarreal, Public domain}}

\end{figure}

Das Skelettsystem besteht aus Knochen und dient 
hauptsächlich als

\begin{itemize}
    \item Stützgerüst  des Körpers
    \item Ursprungs-  und Ansatzfläche für die Muskulatur
\end{itemize}

und weiters als

\begin{itemize}
    \item Kalziumspeicher ($Ca^{++}$)
    \item Blutbildungsorgan (rotes Knochenmark!)
\end{itemize}

\subsection{Der Knochen}
\index{Knochen}

Je nach ihrer Form unterteilt man die Knochen in: 
\index{Knochen!platte}\index{Knochen!Röhren-}\index{Knochen!Würfel-}
\begin{description}
    \item[\gT{Platte Knochen}] (z.B. Beckenknochen,
    Schulterblatt, Brustbein, Schädel),
    \item[\gT{Röhrenknochen}] (Oberarmknochen, Schienbein)
    und
    \item[\gT{Würfelknochen}] (z.B. Hand-, Fußwurzelwurzelknochen,
    Wirbelknochen)
\end{description}


\gAuthNotes{GABS \& KOCH: evt. umstellen auf deutsch -- latein, aber was heißt
Kortikalis auf deutsch?}

\gLayout{24}{Aufbau eines Knochens}{}{

\begin{description}
    \item[\gT{Periost}] \index{Periost}\gE{Beinhaut}, überzieht den Knochen und ist sehr gut mit
Nerven versorgt.  Verantwortlich für Dickenwachstum und Frakturheilung
\item[\gT{Kompakta}] (Kortikalis)
\index{Kompakta}\index{Kortikalis}feste äußere Schicht,
äußere Form
\item[\gT{Spongiosa}] \index{Spongiosa}\index{Knochenbälkchen}Knochenbälkchen ${\rightarrow}$ Zug- und Druckfestigkeit
\item[\gT{Markhöhle}] \index{Markhöhle}s.u.
\item[\gT{Epiphysenfuge}] \index{Epiphysenfuge}Wachstumsfuge
\item[\gT{Gelenkfläche}] \index{Gelenkfläche!Aufbau des
Knochens}mit knorpeligen Überzug als Teil eines Gelenkes.
Manche Knochen haben keine Gelenksfläche.
\end{description}

Lange Röhrenknochen kann man außerdem in 
\begin{compactitem}
    \item \gT{Kopf}\index{Kopf!Knochenaufbau},
    \item \gT{Hals}\index{Hals!Knochenaufbau} und
    \item \gT{Schaft}\index{Schaft!Knochenaufbau}
\end{compactitem}
unterteilen.


}
{./images/anatomie/roehrenknochen-querschnitt.png}
{Ein langer Röhrenknochen im Querschnitt.}
{Lena Hirtler, modifiziert.}

\subsubsection{Das Knochenmark}
\index{Knochenmark}

\paragraph{Das \gT{Rote Knochenmark} ist für die Blutbildung
zuständig}
\index{Knochenmark!rotes}

Das rote Knochenmark ist das blutbildende Organ
schlechthin. Hier werden die \gEs{zellulären Bestandteile des Blutes
gebildet}. Dazu gehören die \gE{Roten Blutkörperchen} (\gE{Erythrozyten}), die
\gE{weißen Blutkörperchen} (\gE{Leukozyten}) sowie die \gE{Blutplättchen}
(\gE{Thrombozyten}) -- \gSee{topic:blut}. 

Bei der Geburt findet man das rote Knochenmark in
sämtlichen Knochen, im Laufe des Lebens bildet es sich zurück und verwandelt sich in {\quotedblbase}Gelbes Fettmark{\textquotedblleft}. Rotes Knochenmark findet man dann lediglich in den platten Knochen wie Becken, Brustbein
und Schädel, wo es für die Blutbildung verantwortlich
ist und normalerweise zeitlebens erhalten bleibt. 

\gAuthNotes{KOCH: wichtig? ev. weglassen?

GABS: Ja, wichtig wegen KT, onko-Patienten, etc. Soll
bleiben. 

KOCH: ok.} 

\gFazit{Einige medizinische Behandlungsmethoden
verursachen als Nebenwirkung eine massive \gEs{Schädigung
des Knochenmarks}\index{Knochenmark!Schädigung} (z.B.
Bestrahlung, Chemotherapien bei Krebspatienten ("`onkologische
Patienten"'), etc.). Diese Patienten leiden neben einer
Blutarmut auch an einem schwachen Immunsystem, da nicht genügend Weiße Blutkörperchen gebildet werden. Sie können daher sehr
leicht schwerwiegende Infektionen bekommen. 

Diese Patienten
sind im Krankentransport häufig anzutreffen.} 

\paragraph{Das \gT{Gelbe Knochenmark} besteht aus Fett} 
\index{Knochenmark!gelbes}

Im Laufe des Lebens bildet sich in den
Röhrenknochen das rote Knochenmark  zu fettreichem, gelben
Knochenmark um.

Das gelbe Fettmark kann bei schweren Knochenverletzungen von Bedeutung
sein, da verletzungs\-bedingt kleine Fetttröpfchen in die
Blutbahn gelangen können, die dann als Fettgerinnsel  zu ernsten
Komplikationen führen können
(Fettembolie\gAuthNotes{+ KOCH}).

\subsubsection{Knochenwachstum}~

Das Skelett eines Neugeborenen besteht bei der Geburt zum größten
Teil aus Knorpel. Erst mit zunehmendem Lebensalter beginnt dieser von
Knochenkernen ausgehend, zu verknöchern. Eine
\gEs{Wachstumszone}
(\index{Epiphysenfuge}\gT{Epiphysenfuge}) bleibt dabei allerdings erhalten.

Die Epiphysenfuge enthält beim Kind/Jugendlichen bis
etwa zum Ende der Pubertät \gE{teilungsfähiges
Knorpelgewebe}, daher ist ein \gEs{Längenwachstum}
möglich. Beim Erwachsenen sind diese knorpeligen Fugen
hingegen vollständig \gE{verknöchert}, sodass
keinerlei Wachstum mehr stattfinden kann.

\subsubsection{Gelenke}
\index{Gelenke}

\gDefinition{(Meistens bewegliche) Verbindung zwischen
mindestens zwei Knochen.}


\gLayout{24}{Prinzipieller Aufbau}{}{
\begin{itemize}
\item Ein Gelenk ist von einer bindegewebigen Kapsel umgeben und durch
Sehnen und Bänder stabil verbunden.

\item Die Gelenksflächen sind zur Reibungs\-ver\-minderung mit 
Knorpel überzogen und der Gelenksspalt mit einer dünnen Schicht 
Schmierflüssigkeit \gZ{(Synovialflüssigkeit)} gefüllt. 

\end{itemize}}
{./images/anatomie/anatomie-gelenk-edited.png}
{\gAuthNotesPicLegalOk{}{Lena }Ein Gelenk
im Querschnitt.}{Lena
Hirtler}\gAuthNotes{KOCH. löschen: Membrana synovialis
und membrana fibrosa}


\paragraph{Es gibt verschiedene Arten von Gelenken}
\index{Gelenke!Arten}

Klassische Beispiele: \index{Kugelgelenk}\gT{Kugelgelenk} (Hüftgelenk), \index{Scharniergelenk}\gT{Scharniergelenk}
(Fingergliedgelenke), \index{Sattelgelenk}\gT{Sattelgelenk}
(Daumengrundgelenk).  


\begin{minipage}{\linewidth}
\includegraphics[width=\linewidth]{./images/noauth/AASdev20090228-img23.png}
\gCaptionof{figure}{\gAuthNotesPicLegalNotOk{img23}{Arten von
Gelenken}{Arten von
Gelenken}}  
\end{minipage}


\subsection{Der Schädel} 
\index{Schädel}
\index{Caput}
\index{Cranium}
\label{topic:schaedel}

Kopf (\gT{Caput}) / Schädel (\gT{Cranium})

 


\begin{gWideParEnv}
\gParallel{
\includegraphics[width=\linewidth]{./images/anatomie/Human_skull_side_simplified_hirngesicht_de-edited.png}
\gCaptionof{figure}{\gAuthNotesPicLegalOk{}{PD}
\index{Hirnschädel}\index{Gesichtsschädel}Der Schädel
gliedert sich in den \gT{Hirnschädel} und den \gT{Gesichtsschädel}}. 
\gSourcePicture{Mariana Ruiz Villarreal, Public Domain}
}
{
\includegraphics[width=\linewidth]{./images/anatomie/Human_skull_front_simplified_bones_de-edited.png}
\gCaptionof{figure}{\gAuthNotesPicLegalOk{}{PD}Schädel Vorderansicht.
\gSourcePicture{Mariana Ruiz Villarreal, Public Domain}}
% \includegraphics[width=\linewidth]{./images/anatomie/Human_skull_side_simplified_bones_de-edited.png}
% \gCaptionof{figure}{\gAuthNotesPicLegalOk{}{PD}Schädel seitlich.
% \gSourcePicture{Mariana Ruiz Villarreal, Public Domain}}
}\end{gWideParEnv}

\gA{STEU 2009-06-28: Genaue Bezeichnungen der einzelnen
Schädelknochen wirklich notwendig? Unterkiefer, Oberkiefer,
Kalotte.

EMHO: Schrift grau. }

\gLayout{0}{}
{
\begin{itemize}
    \item Schädelbasis \index{Schädelbasis}
    \item Schädeldach (Kalotte) \index{Schädeldach}
    \index{Kalotte}
    \item Die platten Knochen sind durch Schädelnähte
    \index{Schädelnähte} verzahnt miteinander verknöchert
\end{itemize}
}
{
\begin{itemize}
    \item Schädelbasis \index{Schädelbasis}
    \item Schädeldach (Kalotte) \index{Schädeldach}
    \index{Kalotte}
    \item Die platten Knochen sind durch Schädelnähte
    \index{Schädelnähte} verzahnt miteinander verknöchert
\end{itemize}
}
{}{}{}




\index{Fontanelle}



\gLayout{21}{Fontanelle}{Die \gT{Fontanellen} sind \gEs{knorpelige 
Verbindungen} zwischen Schädelknochen beim Neu\-geborenen
(der Schädel muss sich für Geburt verformen
können!). \gZ{Sie schließen  sich bis zum 18.
Lebensmonat,} verknöchern  aber erst vollständig mit
zunehmendem  Lebensalter \gZ{(ca. 20 Lebensjahr
abgeschlossen)}.

\begin{itemize}
\item \gZ{große Fontanelle (rot): zw. Stirnbein u.
Scheitelbeinen}
\item \gZ{kleine Fontanelle (grün): zw. Scheitelbeinen u.
Hinter\-haupts\-bein}
\end{itemize}
}{}
{./images/anatomie/fontanelle-dt.jpg}
{\gAuthNotesPicLegalOk{Fontanelle}{Gray}Fontanellen.}
{Gray's Anatomy, Copyright expired.}

% Die \gT{Sutur} ist die endgültig fest verknöcherte Naht
% zwischen den Schädelknochen beim Erwachsenen ${\rightarrow}$ starrer  stabiler 
% Schädel, Schutz  für Gehirn.


\begin{gWideParEnv}
\gLayout{22}{Schädelbasis}{

% a ... Stirnbein (Os frontale)
% 
% b ... Siebbein (Os ethmoidale)
% 
% c ... Keilbein  (Os sphenoidale)
% 
% d ... Schläfenbein(e)  (Os temporale)
% 
% e ... Felsenbein  (Pars petrosa ossis temporalis)
% 
% f ... Hinterhauptsbein  (Os occipitale)
% 
% g ... großes Hinterhauptsloch (Foramen magnum)

Das Gehirn liegt auf der Schädelbasis \index{Schädelbasis}
auf \gZ{(außer wenn man auf dem Kopf steht)}. Hirnnerven treten durch die diversen Löcher aus dem
Schädel aus. \gEs{Das \gT{große
Hinterhauptsloch}\index{Großes
Hinterhauptsloch}  \gZ{(Foramen magnum)} ist die
Durchtrittsstelle des Rückenmarks.} 

\gCave{\index{Hirnstammeinklemmung}\gT{Hirnstammeinklemmung} 

Das \gEs{große Hinterhauptsloch} ist
die Durchtrittsstelle des Rückenmarks. Bei einer
\gEs{Hirndrucksteigerung} (\gE{Hirnödem, Hirnblutung}) ist
es die einzige Möglichkeit zur Ausdehnung des Gehirns!
\gEs{Dabei drückt sich das Gehirn allerdings selber ab},
es kommt zu einer tödlichen \gT{Hirnstammeinklemmung!} ${\rightarrow}$ Lebensgefahr! }}
{}
{./images/anatomie/cavernous_sinus.png}
{}
{Gray's Anatomy, Copyright expired.}


\end{gWideParEnv}

\subsection{Die Wirbelsäule}
\index{Wirbelsäule}
\begin{gWideParEnv}
\gLayout{24}{}{}{
Die Wirbelsäule besteht aus:

\begin{itemize}
    \item Wirbeln \gZ{(vertebrae)}
    \item \index{Bandscheiben}\gT{Bandscheiben}
    (\gT{Discus}): dienen der "`Stoßdämpfung"'
    \item \gEs{Im Wirbelkanal liegt das Rückenmark
    eingebettet.}
\end{itemize}

\vspace{2em}

Sie gliedert sich in:

\begin{itemize}
    \item 7 \gEs{Hals}wirbel \gZ{(Cervical-Wirbel)}
    \index{Halswirbel}
    \item 12 \gEs{Brust}wirbel  \gZ{(Thorakal-Wirbel)}
    \index{Brustwirbel}
    \item 5 \gEs{Lenden}wirbel  \gZ{(Lumbal-Wirbel)}
    \index{Lendenwirbeln}
    \item 5 verschmolzene \gEs{Kreuzbein}wirbel \gZ{(Os
    sacrum)} \index{Kreuzbeinwirbel}
    \item 3-4 \gEs{Steißbein}"`wirbel"' 
    \gZ{(Os coccygis)} \index{Steißbeinwirbel}
\end{itemize}

Die Wirbelsäule hat ein doppelt S-förmig gekrümmtes
Achsenskelett, welches den  aufrechten Gang ermöglicht.
}
{./images/anatomie/wirbelsaeule-lena-edited.png}
{\gAuthNotesPicLegalOk{}{}Die Wirbelsäule.}
{Hirtler.}




\end{gWideParEnv}



\subsubsection{Die Wirbel: Aufbau und Arten}~
\index{Wirbel}

\gLayout{24}{}{}{
Prinzipieller Aufbau:

\begin{itemize}
\item Wirbelkörper \index{Wirbelkörper}
\item Wirbelbogen  \index{Wirbelbogen} (umschließt das
Rückenmark)
\begin{itemize}
\item 2 Querfortsätze \index{Querfortsätze}
\item 1 Dornfortsatz \index{Dornfortsatz}
\end{itemize}
\end{itemize}


In den jeweiligen Abschnitten sehen die Wirbel
unterschiedlich aus.
}{./images/anatomie/Gray94.png}
{Typischer Wirbelkörper.}
{Gray's Anatomy, Copyright expired.}






\subsection{Der Brustkorb
(\index{Brustkorb}\index{Thorax}Thorax)}

Der  Brustkorb besteht vorne aus dem \gT{Brustbein}
\index{Brustbein} (\gT{Sternum}\index{Sternum}) und 12
Rippenpaaren, welche hinten mit der Wirbelsäule in gelenkiger Verbindung stehen. 

\index{Rippen}\index{Rippen!echte}\index{Rippen!unechte}\index{Rippen!fliegende}
Die obersten 7 Rippen-Paare werden als 
{\quotedblbase}\gT{echte Rippen}{\textquotedblleft}
bezeichnet. Sie erreichen mit ihren Spitzen das Brustbein und bilden praktisch
{\quotedblbase}Ringe{\textquotedblleft} um die Brusteingeweide. Die 5
unteren Rippenpaare bilden den bogenförmigen unteren Rand des
Brustkorbes. Die Rippenpaare 8-10 sind dabei nur indirekt
über einen Knorpel mit dem Brustbein verbunden
({\quotedblbase}\gT{unechte Rippen}{\textquotedblleft}), das
11. und das 12. Paar endet mit knorpeligen Spitzen frei zwischen den
Bauchmuskeln({\quotedblbase}\gT{fliegende
Rippen}{\textquotedblleft}) bezeichnet werden.

Der Brustkorb enthält und schützt die Brust\-organe (Lunge, Herz)
sowie die Oberbauchorgane (Leber, Magen). Zwischen den
Rippen spannt sich die
\gT{Zwischenrippenmuskulatur}\index{Zwischenrippenmuskulatur} (wichtig für \gEs{Atemmechanik}), in diese sind segmentale
Nerven- und Gefäßbündel eingelagert. Das \gEs{Zwerchfell}
\index{Zwerchfell} schließt den Brustraum nach unten hin ab
und trennt somit den Brustraum vom Bauchraum.

\gLayout{23}{}
{
Der Thorax besteht aus:

\begin{itemize}
    \item \gEs{Brustbein} (\gT{Sternum})
    \item \gEs{12  Rippenpaaren}
    \begin{itemize}
        \item \gZ{7 "`echte"'}
        \item \gZ{3 "`unechte"'}
        \item \gZ{2 "`fliegende"'}
    \end{itemize}
    \item \gEs{Brustwirbelsäule} (12 Brust\-wirbel), sie
    steht mit den Rippen in gelenkiger Verbindung
\end{itemize}
(\emph{{\quotedblbase}echt{\textquotedblleft}} bedeutet die
Rippe hat eine direkte Verbindung zum Sternum)
}{}{./images/anatomie/thorax.png}{Der knöcherne Thorax.}{Mariana Ruiz Villarreal, Public Domain}



\subsection{Der \index{Schultergürtel}Schultergürtel
und die obere Extremität}
\index{Schultergürtel}
\label{topic:schulterguertel}


\begin{figure}[H]
    \begin{center}
    \includegraphics[width=\linewidth]{./images/anatomie/Human_arm_bones_diagram_de-edited.png}
\gCaptionof{figure}{\gAuthNotesPicLegalOk{}{PD}Der
Schultergürtel und die obere Extremität.
\gSourcePicture{Mariana Ruiz Villarreal, Public Domain}}
    \end{center}
\end{figure}

\gParallel{\subsubsection{Der Schultergürtel} 

\ldots besteht aus (von vorne nach
hinten):

\begin{itemize}
    \item \index{Brustbein}\gT{Brustbein}
    (\index{Sternum}\gT{Sternum})
    \item \index{Schlüsselbein}\gT{Schlüsselbein}
    \gZ{(Clavicula)}
    \item \index{Schulterblatt}\gT{Schulterblatt} \gZ{(Scapula)}
\end{itemize}

Der Schultergürtel bildet die Gelenkspfanne für das Schultergelenk,
jeweils auf beiden Seiten.

Die Schulterblätter sind nur durch Muskulatur mit dem Brustkorb
verbunden.
}{
\subsubsection{Die obere Extremität}
\index{Extremität!obere}
\label{topic:obere-extremitaet}

\begin{itemize}
    \item \gEs{Schultergelenk} \index{Schultergelenk}
    \item \gEs{Oberarmknochen} \index{Oberarmknochen} 
    (\gT{Humerus}\index{Humerus})
    \item \gEs{Ellenbogengelenk} \index{Ellenbogengelenk}
    \item \gEs{Unterarmknochen} \index{Unterarmknochen} (2)
    \begin{itemize}
        \item \gEs{Elle} \index{Elle}
        (\gT{Ulna}\index{Ulna})
        \item \gEs{Speiche} \index{Speiche}
        (\gT{Radius}\index{Radius}), daumenseitig!
    \end{itemize}
    \item Handwurzelgelenk \index{Handwurzelgelenk}
    \item (8) \gEs{Handwurzelknochen}
    \index{Handwurzelknochen}
    \item (5) \gEs{Mittelhandknochen}
    \index{Mittelhandknochen}
    \item \gEs{Fingergliedknochen}
    \index{Fingergliedknochen} mit 2-3 Gliedern
\end{itemize}
}






\gAuthNotes{GABS: Bild evt. nur mit einer Seite, dafür aber
größer und neben Text?

KOCH: passt so.}


\gFazit{Achtung: Unterscheide: Hand --
Arm!\index{Hand}\index{Arm}}









\subsection{Der Beckengürtel und die untere Extremität}
\index{Beckengürtel}
\index{Extremität, untere}
\label{topic:beckenguertel}









\gLayout{23}{}
{
\subparagraph{Der Beckengürtel}
Besteht aus (von hinten nach vorne) :

\begin{itemize}
    \item Kreuzbein (Os sacrum)
    \item Beckenknochen, bestehend aus:
    \begin{itemize}
        \item \gZ{\index{Darmbein|gZ}Darmbein}
        \item \gZ{\index{Sitzbein|gZ}Sitzbein}
        \item \gZ{\index{Schambein|gZ}Schambein}
    \end{itemize}
\end{itemize}
\gZ{Die \index{Symphyse|gZ}Symphyse stellt eine straffe
Knorpelverbindung vorne zwischen den beiden Beckenknochen
dar (Dehnungs\-fuge für Geburt).

\subparagraph{Die untere Extremität}
\label{topic:untere-extremitaet}

\begin{itemize}
    \item \gEs{Hüftgelenk} \index{Hüftgelenk} 
    \item \gEs{Oberschenkel} \index{Oberschenkel}  
    (\gT{Femur})
    \index{Femur} 
    \item \gEs{Kniegelenk} \index{Kniegelenk} 
    \item \gEs{Kniescheibe} \index{Kniescheibe} 
    \item \gEs{Unterschenkelknochen} (2)
    \index{Unterschenkelknochen}

    \begin{itemize}
        \item \gEs{Schienbein}  (Tibia) , dick, vorne
        \index{Schienbein}  \index{Tibia} 
        \item \gEs{Wadenbein}  (Fibula), dünn, seitlich
        hinten \index{Wadenbein}  \index{Fibula} 
    \end{itemize}
    \item \gEs{Sprunggelenk} \index{Sprunggelenk} 
    \item 7 \gEs{Fußwurzelknochen}, \gZ{u.a. Sprungbein,
    Fersenbein} \index{Fußwurzelknochen} 
    \item 5 \gEs{Mittelfußknochen}  
    \index{Mittelfußknochen}
    \item \gEs{Zehengliederknochen} mit je 2-3 Gliedern
    \index{Zehengliederknochen} 
\end{itemize}}
}{}{./images/anatomie/untere-extremitaet.png}{\gAuthNotesPicLegalOk{}{}Der
Beckengürtel und die untere Extremität.}{Mariana Ruiz Villarreal, Public Domain}


% TODO INHALT: Knöchel erwähnen (STEU)
\gA{STEU: Knöchel erwähnen.}

\gFazit{Achtung: Unterscheide Fuß -- Bein!}

\subsection{Die Muskulatur}
 \index{Muskel}  \index{Muskulatur!glatte} 
 \index{Muskulatur!quergestreifte}
 \index{Muskulatur!Skelett-}
 \index{Muskulatur!Herz-}

Prinzipiell gibt es 3  Arten von Muskelzellen:

\begin{itemize}
    \item \gT{Quergestreifte (Skelett-)Muskulatur}:
    willkürlich steuerbar, trainierbar, er\-müd\-bar
    \item \gT{Glatte  Muskulatur}: Nicht willentlich
    kontrollierbar, unermüdlich, ausdauernd
    \item \gT{Herzmuskel}: Sonderstellung zwischen glatt und
    quergestreift, Eigenschaften von beiden: nicht willentlich
    steuerbar, trainierbar, nicht ermüdend, quergestreift.
\end{itemize}

Die \gE{quergestreifte Muskulatur} heißt deswegen
{\quotedblbase}quergestreift{\textquotedblleft}, da man unter dem
Mikroskop tatsächlich eine Streifung sehen kann, allerdings nicht mit
dem freien Auge. 

\gLayout{22}{Gegenspieler}
{
Muskel können nur Zug und keinen
Druck ausüben, daher sind \gEs{Gegenspieler} (Agonist u.
Antagonist) erforderlich (z.B. ein Muskel streckt den Arm,
ein \gE{anderer} beugt ihn). Jeder Muskel hat einen Ursprung
und einen Ansatz, diese beiden Punkte bestimmen die Bewegung.
}{}{./images/anatomie/Gray361.png}{\gAuthNotesPicLegalOk{}{}Muskelzug. }{Gray's
Anatomy, Copyright expired.}






%%%%%%%%%%%%%%%%%%%%%%%%%%%%%%%%%%%%%%%%%%%%%%%%%%%%%%%%%%%
%%%%%%%%%%%%%%%%%%%%%%%%%%%%%%%%%%%%%%%%%%%%%%%%%%%%%%%%%%%
%%%%%%%%%%%%%%%%%%%%%%%%%%%%%%%%%%%%%%%%%%%%%%%%%%%%%%%%%%%
%%%%%%%%%%%%%%%%%%%%%%%%%%%%%%%%%%%%%%%%%%%%%%%%%%%%%%%%%%%
\clearpage
\section{Das Herz}
 \index{Herz} 

\gLayout{22}{}
{
Das Herz ist ein etwa faustgroßer \gEs{Hohlmuskel} und
hat die Funktion einer \gEs{Pumpe}. Es liegt \gEs{hinter dem
Sternum\footnote{Brustbein}}\gA{\footnote{\gA{EMHO: statt
Fußzeile: "`Brustbein (Sternum)"'. damit einheitlich.\par
GABS: nein. Hier sollen sie bereits primär den Begriff
"`Sternum"' verwenden udn daran gewöhnt werden.}}} und
teilweise im linken Brustkorb. Unten liegt es auf \gEs{dem Zwerchfell auf}.

 \index{Scheidewand}  \index{Septum}  
 \index{Vorhof}
 \index{Atrium}  \index{Herzkammer}
 \index{Ventrikel}  
 Das Herz ist durch eine \gEs{Scheidewand} (Septum) in eine
\gE{rechte} und eine \gE{linke} Seite geteilt.
Umgangssprachlich spricht man dabei jeweils vom
\emph{{\quotedblbase}rechten{\textquotedblleft}} oder
\emph{{\quotedblbase}linken Herz{\textquotedblleft}}. Jede
Hälfte wird durch \gEs{Klappen} in einen \gT{Vorhof}
(\gT{Atrium}) und eine \gT{Herzkammer} (\gT{Ventrikel})
geteilt.

\subparagraph{Blutversorgung}

Die Blutversorgung erfolgt über die
\gT{Herzkranzgefäße} \index{Herzkranzgefäße} 
(\gT{Koronargefäße} \index{Koronargefäße} ) welche direkt
aus der Hauptschlagader abgehen.

\subparagraph{Schichten}

3 Schichten (innen nach aussen):
\gAuthNotes{\footnote{GABS: Schichten relevant? KOCH: ja!}}
\begin{itemize}
\item Endokard (innen) \index{Endokard}     
\item \gEs{Myokard} (\gEs{Muskel}) \index{Myokard|gEs} 
\item Epikard (aussen) \index{Epikard} 
\item Perikard (ganz aussen), bildet den Herzbeutel.
\end{itemize}
}{}{./images/noauth/AASdev20090228-img27.png}{\gAuthNotesPicLegalNotOk{img27}{}Das Herz.}{Quelle nicht eruierbar.}

~


\gLayout{22}{Herzklappen}
{
\index{Herzklappen}   \index{Segelklappen}
 \index{Taschenklappen}

Insgesamt \gEs{4
Herzklappen} sorgen als Ventile dafür dass das Blut nur in eine Richtung strömt. Es gibt 2 Segelklappen und 2 Taschenklappen.
Die Segelklappen befinden sich zwischen den Vorhöfen und den
Herzkammern, die Taschenklappen befinden sich an den
Ausflussöffnungen der Kammern am Übergang zur Lungenarterie (rechte
Seite) bzw. zur  Aorta (Hauptschlagader, linke Seite). 
}
{}
{./images/noauth/AASdev20090228-img28.jpg}
{\gAuthNotesPicLegalNotOk{img28}{}Das
 Herz.\gAuthNotes{G-FDL !!!, KOCH: Bild gefällt mir nicht. }}
 {Wikipedia,  Jörg Rittmeister. G-FDL, CCAL}







 
% \gCaptionof{figure}{ Schema der
% Herz-Anatomie.  }


\paragraph{Herzfunktion}\label{topic:herzfunktion}

Der Herzmuskel pumpt und entspannt sich beim Erwachsenen
in Ruhe ca. 60 bis 100~mal in der Minute. Bei jeder
Anspannung (Kontraktion) wirft das Herz ca. 70\,ml Blut in die Arterien
aus. Dabei pumpt es ca. 4\,-\,6\,l  Blut (bzw. 1\,x das
gesamte Blutvolumen) in der Minute.

Durch die Pumpfunktion des Herzens entsteht ein
lebenswichtiger Druck im (arteriellen) Gefäßsystem, der
(arterielle systolische) \gE{Blutdruck} \index{Blutdruck} .
Vgl. Kap. \emph{Vitalfunktionen}, \ref{topic:blutdruck}. 

%
%
\subsection{Reizleitungssystem}
 \index{Reizleitungssystem} 

\gAuthNotes{GABS @ STEU: Reizleitungssystem -- wie
ausführlich? Prüfungsfrage?

KOCH: Beim Defi schon.

GABS: Defi-Teil noch nicht abgefasst!}

\gLayout{23}{}
{
Eigenständiges  Erregungs- und  Reizleitungssystem, streng
hierarchisch, aber jede  Station zur Erregungsbildung  fähig 

bestehend aus:

\begin{itemize}
\item Sinusknoten  \index{Sinusknoten} 
({\sffamily 1}):  
{\quotedblbase}Schrittmacher{\textquotedblleft}, feuert mit einer Eigenfrequenz v. ca. 60/min. elektrische Impulse ab)
\item AV-(AtrioVentricular)-Knoten ({\sffamily 2}): Leitung
vom Vorhof in die Kammern, filtert zu hohe Frequenzen
\index{AV-Knoten} 
\item His-Bündel \index{His-Bündel} 
\item \gZ{Tawara-}Schenkel (linker/rechter)
\item \gZ{Purkinje-}Fasern
({\quotedblbase}Endstrecke{\textquotedblleft})
\end{itemize}

\gFazit{Der normale Schrittmacher ist der Sinusknoten,
welcher die elektrischen Impulse für die Vorhöfe und die Kammern erzeugt. Man spricht dann
vom Sinusrhythmus.}

Beispiele von typischen Rhythmen werden im Kapitel
Vitalfunktionen vorgestellt \gAuthNotes{Referenz einfügen!}.
}{}{./images/anatomie/Reizleitungssystem_1.png}{\gAuthNotesPicLegalOk{}{CC-BY}
Reizleitungssystem, Schema.}{J. Heuser,
basierend auf der Arbeit von Patrick J. Lynch; illustrator;
C. Carl Jaffe; MD; cardiologist Yale University Center for
Advanced Instructional Media. Lizenz: CC-BY}






\gAuthNotes{KOCH: Tabelle mit Normalwerten einfügen

GABS: in Kap Vitalfunktionen}


\section{Der Kreislauf }
 \index{Kreislauf}
\label{topic:blutkreislauf} 

\gFazit{\gEs{Herz + Gefäßsystem + Blut = Kreislauf}}

Der Blutdruck und die Pumpfunktion des Herzens werden im
Kapitel \emph{Vitalfunktionen} unter \ref{topic:blutdruck}
(Blutdruck) sowie im Abschnitt \ref{topic:herzfunktion}
(Herzfunktion) in diesem Kapitel besprochen.

\index{Körperkreislauf}  \index{Kreislauf!großer} 
\index{Kreislauf!kleiner}
\index{Lungenkreislauf}
\index{Hochdrucksystem}  \index{Niederdrucksystem}
Es wird der \gT{Körperkreislauf} (\gE{"`Großer
Kreislauf"'}, "`Hochdrucksystem"') vom
\gT{Lungenkreislauf} (\gE{"`Kleiner Kreislauf"'},
"`Niederdrucksystem"') unterschieden. Das Herz ist für 
beide Systeme die gemeinsame Pumpe.

Der Kreislauf wird vom \gEs{Stammhirn} gesteuert
(Herzfrequenz, Blutdruckregulation, etc.).

\subsection{Abschnitte des Kreislaufs}
\vspace{-1em} \index{Gasaustausch} 
\index{Kreislauf!Abschnitte}

\begin{gWideParEnv}
\gParallel
{\subsubsection{Kleiner (Lungen-)Kreislauf}
Die obere und untere Hohlvene münden in den \gEs{rechten
Vorhof} des Herzens und liefern sauerstoffarmes Blut aus
dem Körper an. Von hier wird das Blut in die \gEs{rechte
Herzkammer} gepumpt, um dann in die \gEs{Lungenarterie} zu
gelangen. In der Lunge verästeln sich die Lungenarterien
zu haardünnen Gefäßen, den \gEs{Kapillaren}. Diese
umspannen die Lungenbläschen es findet ein
\gEs{Gasaustausch} statt: Sauerstoff wird aufgenommen und
Kohlendioxid abgegeben. Das jetzt mit Sauerstoff
angereicherte Blut wird über die \gEs{Lungenvenen} in den
\gEs{linken Vorhof} geleitet, der wiederum das Blut in die
\gEs{linke Herzkammer} pumpt.

\vspace{1em}

\subsubsection{Großer Kreislauf}
Das Blut wird aus der linken Herzkammer ausgestoßen und
gelangt über die die Hauptschlagader
("`\gEs{Aorta}"') in den
Körperkreislauf. Aus der Aorta entspringen \gEs{Arterien}
({\quotedblbase}Schlagadern{\textquotedblleft}), die sich
immer weiter verzweigen und an Durchmesser abnehmen. Sie
versorgen den gesamten Körper mit Blut. Der durch die
Pumpleistung des Herzens erzeugte Puls ist an oberflächlich
gelegenen Arterien tastbar (Radialis-Puls, Carotis-Puls, ..).
Daher rührt auch der Name
"`Schlagader"'.

Sie verzweigen sich, weiter an Durchmesser abnehmend,
ähnlich dem Geäst eines Baumes, bis hin zu
\gEs{Arteriolen} und haardünnen Gefäßen, den
Haargefäßen ("`\gEs{Kapillaren}"').
Hier findet wieder ein \gEs{Gasaustausch} statt, diesmal
allerdings spiegelbildlich zum Lungenkreislauf:
\gEs{Sauerstoff wird an das Gewebe abgegeben, im Gegenzug
wird Kohlendioxid abtransportiert}.

Das nun sauerstoffarme Blut gelangt über kleine
\gEs{Venolen} in das Venensystem und fließt in den
\gEs{Venen} herzwärts. Am Weg zum Herzen nehmen die Venen
an Umfang immer mehr zu, bis sie letztendlich in die großen Venen
({\quotedblbase}\gEs{obere} und \gEs{untere
Hohlvene}{\textquotedblleft}) münden. 
}
{
\includegraphics[width=\linewidth]{./images/noauth/kreislauf-schema.png}

\gCaptionof{figure}{\gAuthNotesPicLegalNotOk{}{Blutkreislauf}Der
Blutkreislauf.}

\gFazit{Rechter Vorhof ${\rightarrow}$ rechte Herzkammer
${\rightarrow}$ Lungenarterie(n) ${\rightarrow}$
Lungenkapillaren (\ce{CO2}-Abgabe, \ce{O2}-Aufnahme) ${\rightarrow}$ Lungenvenen ${\rightarrow}$ linker Vorhof 

${\rightarrow}$ linke Herzkammer ${\rightarrow}$ Aorta ${\rightarrow}$
abzweigende Arterien ${\rightarrow}$ Kapillaren
(\ce{O2}-Abgabe, \ce{CO2}-Aufnahme)  ${\rightarrow}$ Venen
${\rightarrow}$ Hohlvene ${\rightarrow}$ Rechter Vorhof

 ${\rightarrow}$ alles beginnt wieder von vorne
 $\circlearrowright$}}
\end{gWideParEnv}

~




\subsection{Die Blutgefäße}

Man unterscheidet zwischen:

\begin{description}
    \item[\gT{Arterien}] führen vom Herzen weg, dicke
    Muskelschicht. \index{Arterien}
    \item[\gT{Venen}] führen \gE{zum Herzen hin},
    besitzen Venen\gE{klappen} um Rückfluss\footnote{Rückfluss:
    Fluss entgegen der Flussrichtung} zu vermeiden.
    \index{Venen}
    \item[\gT{Kapillargefäße}] dünne "`Haargefäße"'
    in denen der Gasaustausch stattfindet, Verbindung zwischen Arterien
    und Venen. \index{Kapillargefäße}
\end{description}

Arterien besitzen eine dicke Muskelschicht, Venen
hingegen haben eine deutlich dünnere Wand als die
muskulösen Arterien. Ein \gT{Shunt} \index{Shunt}  ist eine
künstlich (chirurgisch) angelegte Verbindung zwischen Arterien und
Venen und wird oft bei Dialyse-Patienten 
\index{Dialyse-Patient!Shunt} eingesetzt (vgl.
{topic:blutdruckmessung}).

\subsubsection{Wichtige Blutgefäße}

\begin{gWideParEnv}

{\sffamily
\label{topic:blutgefaesse-wichtige}

% \begin{minipage}{\linewidth}
\gCaptionof{table}{Übersicht wichtiger Blutgefäße}

\begin{tabulary}{\linewidth}{lllL}
\\\addlinespace
\noindent\gTabHeading{Gefäß}
&
\noindent\gTabHeading{Fachbegriff}
&
\noindent\gTabHeading{Typ}
&
\noindent\gTabHeading{Bemerkungen}
\\\midrule\addlinespace
\gEs{Hauptschlagader} \index{Hauptschlagader}
&
\gT{Aorta} \index{Aorta}
&
Arterie
&
Geht von der linken Herzkammer ab, von dort zweigen alle anderen
Körperarterien ab
\\\addlinespace
\gEs{Halsschlagader} \index{Halsschlagader}
&
\gT{Karotis} \index{Karotis}
&
Arterie
&
1 Paar (2 Stk.), versorgen den Schädel, Puls tastbar seitlich vom
Kehlkopf
\\\addlinespace
\gEs{Herzkranzgefäße} \index{Herzkranzgefäße}
&
\gT{Koronararterien} \index{Koronararterien}
&
Arterie
&
Versorgen das Herz mit Blut, gehen aus der Aorta ab
\\\addlinespace
\gEs{Lungenarterie} \index{Lungenarterie}
&
\gT{Pulmonalarterie} \index{Pulmonalarterie}
&
Arterie
&
\noindent\gEs{\ce{O2}-arm.} Geht von der rechten Herzkammer weg in
die Lunge zu den Lungenkapillaren
\\\addlinespace
\gEs{Oberarmarterie} \index{Oberarmarterie}
&
Brachialarterie
&
Arterie
&
Taststelle zum Puls-fühlen beim Baby, Blutdruckmessung
\\\addlinespace
\gEs{Radialarterie} \index{Radialarterie}
&
\gT{A. radialis} \index{A. radialis}
&
Arterie
&
Verläuft am Unterarm speichenseitig (dort wo der Daumen ist). Puls
fühlbar.
\\\addlinespace
\gEs{Obere Hohlvene} \index{Obere Hohlvene}
\index{Hohlvene!obere}
&
\gT{Vena cava} sup. \index{Vena cava}
&
Vene
&
Obere Vene die in den rechten Vorhof mündet
\\\addlinespace
\gEs{Untere Hohlvene}
\index{Untere Hohlvene}
\index{Hohlvene!untere}
&
\gT{Vena cava} inf. \index{Vena cava}
&
Vene
&
Untere Vene die in den rechten Vorhof mündet.
\\\addlinespace
\gEs{Lungenvene} \index{Lungenvene}
&
\gT{Pulmonalvene} \index{Pulmonalvene}
&
Vene
&
\gEs{\noindent\ce{O2}-reich.} Kommt von den Lungenkapillaren und
mündet in den linken Vorhof
\\\addlinespace
\gEs{Pfortader} \index{Pfortader}
&
\gZ{Vena portae}
&
Vene
&
Führt \ce{O2}-armes, aber nährstoffreiches Blut von Teilen des Darms zur
Leber (zur Entgiftung)
\\\addlinespace
\gEs{Haargefäße} \index{Haargefäße}
&
\gT{Kapillargefäße} \index{Kapillargefäße}
&
Kapillare
&
Dünne, dicht verzweigte Gefäße, Verbindung zwischen Arterien
und Venen. Hier findet der Gas- und Nährstoffaustausch statt. Es gibt
Körperkapillaren (Körperkreislauf) und Lungenkapillaren
(Lungenkreislauf)
\\\addlinespace\bottomrule
\end{tabulary}
% \end{minipage}
}
% STEU: Leistenarterie! Puls
% GABS: Wozu?
\end{gWideParEnv}

\subsection{Das Blut}\index{Blut}\label{topic:blut}

Das Blut macht ca. \gEs{7-8\,\% des Körpergewichtes} aus,
das entspricht in etwa \gEs{5-6\,l}.

Es setzt sich zusammen aus:

\begin{itemize}
    \item festen Bestandteilen (Zellen, Blutkörperchen)
    \item flüssigen Bestandteilen (\index{Plasma}Plasma)
\end{itemize}

\begin{minipage}{\linewidth}
 \includegraphics[width=\linewidth]{./images/noauth/AASdev20090228-img32.png} 
 \gCaptionof{figure}{\gAuthNotesPicLegalNotOk{img32}{}Die
 Bestandteile des Blutes.
 \gSourcePicture{Quelle nicht eruierbar.}}
\end{minipage}


\paragraph{Blut - Aufgaben}

\begin{itemize}
    \item Sauerstofftransport zu den Körperzellen
    \index{Sauerstofftransport}
    \item \ce{CO2}-Abtransport \index{CO2-Abtransport@\ce{CO2}-Abtransport} 
    \item Transport von Nährstoffen und Abfallprodukten
    \item Wärmeverteilung
    \item \gZ{pH-Pufferungssystem (Säure/Basen-Gleichgewicht)
    (pH-Wert muss konstant zwischen 7,35 und 7,45 liegen!)}
\end{itemize}


\subsubsection{Feste (zelluläre) Bestandteile}

\begin{itemize}
    \item \gT{Rote Blutkörperchen} (\gT{Erythrozyten})
    \item \gT{Weiße Blutkörperchen} (\gT{Leukozyten})
    \item \gT{Blutplättchen} (\gT{Thrombozyten})
\end{itemize}

\paragraph{Erythrozyten}\index{Erythrozyten} \index{Rote
Blutkörperchen}
\index{Blutkörperchen!rote} 

\begin{itemize}
    \item Rote Blutkörperchen
    \item Enthalten den roten, eisenhaltigen Blutfarbstoff
    \gT{Hämoglobin}\index{Hämoglobin}, das Hämoglobin ist
    das Transportprotein für Sauerstoff
    \item \ce{O2}-Transport
    \item Tragen Blutgruppenmerkmale (Blutgruppe)
    \item seeeeeehr viele, die häufigsten Blutkörperchen
    \gZ{(4.000.000\,/$\mu$l)}
\end{itemize}

\paragraph{Leukozyten}\index{Leukozyten} \index{Weiße
Blutkörperchen}  \index{Blutkörperchen!weiße}

\begin{itemize}
    \item Weiße Blutkörperchen
    \item Körperabwehr, Entzündungszellen, Immunsystem 
    \index{Immunsystem!Leukozyten} 
    \item
    \gZ{(4\,000\,-\,8\,000\,/$\mu$l)}
\end{itemize}

\paragraph{Thrombozyten}\index{Thrombozyten}

\begin{itemize}
    \item \index{Blutplättchen}Blutplättchen
    \item Wichtige Rolle beim Wundverschluss
    \item Zelluläre Blutstillung
    \item
    \gZ{(150\,000\,-\,300\,000\,/$\mu$l)}
\end{itemize}

\subsubsection{Die flüssigen Bestandteile --
Das Plasma}\index{Plasma}

\begin{itemize}
    \item 10\% Proteine, Enzyme, Hormone, Elektrolyte, Glucose, Fette
    \item 90\% Wasser
    \item Aufgaben:
    \begin{itemize}
        \item Nährstofftransport
        \item Wärmeverteilung
    \end{itemize}
    \item enthält Proteine und Enzyme für
    \gEs{Plasmatische Blutgerinnung} (Serum)
    \item Pufferung des pH-Wertes
\end{itemize}

\subsection{Die Blutstillung}
 \index{Blutstillung} 

\gMarried{
Die Blutstillung umfasst alle Vorgänge, die zwischen dem Entstehen und dem
Verschluss einer Wunde ablaufen. Sie läuft in zwei Phasen ab: vorläufiger
Wundverschluss durch Bildung eines \gE{Thrombozytenpfropfs} 
\index{Thrombozytenpfropf} und endgültiger Wundverschluss
durch Blutgerinnung \index{Blutgerinnung}. 
}{
\includegraphics[width=\linewidth]{./images/anatomie/Bleeding_finger.jpg}
\gCaptionof{figure}{\gAuthNotesPicLegalOk{}{CC-BY}\gSourcePicture{Crystal
(Crystl) from Bloomington, USA
[\url{http://www.flickr.com/people/crystalflickr/}]
Lizenz: CC-BY}} }


\paragraph{Bildung eines Thrombozytenpfropfs} Zunächst
lagern sich an die verletzte und defekte Stelle des Blutgefäßes Thrombozyten an
und verkleben miteinander. Die verklebenden
Blutplättchen bilden einen Thrombozytenpfropf.

\gFazit{Die Bildung des Thrombozytenpfropf wird von den zelluären Bestandteilen
des Blutes erledigt.}

\paragraph{Blutgerinnung} \gZ{Die Blutgerinnung ist der wichtigste
Prozess bei der Blutstillung. Sie wird entweder durch die
Gewebeverletzung selber oder durch ein eigenes System von Gerinnungsfaktoren
ausgelöst. Der Vorgang der Blutgerinnung hat zum Ziel,
das im Blutplasma gelöste \gZ{Fibrinogen} zu aktivieren und
in einen gallertartigen, klebrigen und später festen Zustand
zu überführen. Um dieses Ziel zu erreichen, steht dem Organismus ein Gerinnungsystem bestehend aus im Plasma gelösten Gerinnungsfaktoren zur
Verfügung, das kaskadenartig aufgebaut ist. Dieses System funktioniert wie eine
Kette aufgestellter Dominosteine, die nacheinander umfallen, wenn der erste Stein
angestoßen wird. }

\gFazit{Die Blutgerinnung kommt durch im Plasma gelöste Gerinnungsfaktoren
zustande.}

Neben dem System der Blutgerinnung
gibt es auch ein System zur Auflösung von Thromben, die Fibrinolyse. Beide
Systeme müssen sich die Waage halten.

\gFazit{An der Blutstillung sind sowohl die festen Bestandteile des Blutes als
auch das Plasma beteiligt.}

\clearpage
\section{Das Nervensystem}
 \index{Nervensystem} 

\paragraph{Unterteilung}~

Man kann das Nervensystem nach seinem Aufbau und nach seiner Funktion
unterteilen.

Vom Blickwinkel des \underline{Aufbaus} wird das Nervensystem grob
unterteilt in ein


\begin{itemize}
\item \gT{Zentrales Nervensystem} (\gT{ZNS}) und ein 
\item \gT{Peripheres Nervensystem}.
\end{itemize}

Das ZNS ist im wesentlichen für Steuer-, Regel- und Denkfunktionen
zuständig sowie Teilweise für die Weiterleitung von Impulsen und
Informationen. Das periphere Nervensystem ist fast ausschließlich nur
für die Weiterleitung von Informationen und Impulsen zu den
Endorganen (oder umgekehrt, von den Endorganen zum ZNS) zuständig. 

\marginlabel{Die Begriffe sind recht unscharf
definiert und geben nur eine grobe, aber wichtige Übersicht
wieder.}Betrachtet man die \underline{Funktion} kann man das
Nervensystem u.a. in ein

\begin{itemize}
\item \gT{Motorisches} (willkürliche\footnote{\emph{Willkür}: Dem \gE{freien
Willen} unterworfen.} Bewegung),
\item \gT{Sensorisches} ("`fühlend"', Reizweiterleitung von
den Sinnesorganen) und ein
\item \gT{Vegetatives} (autonomes, bzw. unwillkürliches)
Nervensystem
\end{itemize}

unterteilen. 

\subsection{Das Zentralnervensystem -- ZNS}
 \index{Zentralnervensystem}  
 \index{ZNS} 
\index{Nervensystem!Zentral-}
\label{topic:zns}
\gTagNo{Unterteilung nach Lage und Aufbau}

\gMarried{
Das Zentralnervensystem (ZNS) besteht aus:

\begin{itemize}
    \item Gehirn
    \begin{itemize}
        \item Großhirn
        \item Zwischenhirn
        \item Kleinhirn
        \item Hirnstamm
    \end{itemize}
    \item Rückenmark
\end{itemize}


\gA{STEU: Hirnteile, Balken in grau}

\paragraph{Aufgaben allgemein}~

\begin{itemize}
    \item Aufnahme  und Verarbeitung  von äußeren und internen  Reizen
    (Schmerz, Fühlen, Sehen, Hören, {\dots})
    \item Bewusstsein und {\quotedblbase}Bewusstmachung{\textquotedblleft}
    \item Abgabe  von zweckentsprechenden Impulsen (Steuerung, Reaktion
    etc.)
\end{itemize}
}{
\includegraphics[width=\linewidth]{./images/anatomie/Brain_human_sagittal_section.png}
\gCaptionof{figure}{Das Hirn im seitlichen Querschnitt.
\gSourcePicture{Patrick J. Lynch, medical illustrator,
\url{http://patricklynch.net}; C. Carl Jaffe, MD,
cardiologist. Lizenz: CC-BY 2.5}} }


\paragraph{Spezielle Aufgaben}~

\index{Großhirn}\index{Zwischenhirn}\index{Kleinhirn}\index{Hirnstamm}
\index{verlängertes Rückenmark}\index{Rückenmark}

\gT{Großhirn}

\begin{itemize}
    \item Bewusstes Denken
    \item Sitz des Bewusstseins
    \item Gefühle
    \item Gedächtnis
\end{itemize}

% \gT{Zwischenhirn}
% \begin{itemize}
% \item Schaltinstanz und Filter Rückenmark - Großhirn
% \end{itemize}

\gT{Kleinhirn}

\begin{itemize}
    \item Bewegungskoordination ({\quotedblbase}Makros{\textquotedblleft})
\end{itemize}

\gT{Hirnstamm} und \gT{verlängertes Rückenmark}

\begin{itemize}
    \item Atem-, Kreislaufzentrum
\end{itemize}

\gT{Rückenmark}

\begin{itemize}
    \item Weiterleitung von Impulsen (Leitungsbahnen)
    \item Vom Rückenmark gehen die peripheren Nerven aus
    \item (Muskeleigen-)Reflexe
\end{itemize}


\gMarried{\paragraph{Der Schädel besteht aus Schichten}~
%\index{Epiduralraum}
%\index{Dura mater}
\index{Harte Hirnhaut}
%\index{Subduralraum}
%\index{Arachnoidea}
%\index{Subarachnoidalraum}
\gAuthNotes{GABS: Wie sehr hervorgehoben?}
\begin{itemize}
    \item Kopfschwarte
    \item Knochen
%    \item Epiduralraum
    \item Harte Hirnhaut \gZ{(Dura mater)}
%    \item Subduralraum
    \item Arachnoidea (Spinngewebshaut)
%    \item Subarachnoidalraum
%    \item eingebettete Gefäße
    \item Weiche Hirnhaut \gZ{(Pia mater)}
\end{itemize}}
{
\includegraphics[width=\linewidth]{./images/anatomie/Gray1196-edited.png}
\gCaptionof{figure}{\gAuthNotesPicLegalOk{}{PD-CE}Querschnitt durch
die Schädeldecke. \gSourcePicture{Gray's Anatomy, Copyright
expired.}}}


% Aufbau ZNS
% 
% \begin{itemize}
% \item Großhirn (Telencephalon)
% 
% \end{itemize}
% \begin{itemize}
% \item \index{Graue Substanz}Graue Substanz (Rinde)
% 
% \item \index{Weiße Substanz}Weiße Substanz (Mark)
% 
% \item Balken (Verbindung li/re)
% 
% \end{itemize}
% \begin{itemize}
% \item Zwischenhirn  (Diencephalon)
% 
% \end{itemize}
% \begin{itemize}
% \item Thalamus, Hypothalamus (Hormone, Temperaturregulation etc.)
% 
% \item Hypophyse (Hormone)
% 
% \end{itemize}
% \begin{itemize}
% \item Hirnstamm (Truncus encephali)
% 
% \item Mittelhirn (Mesencephalon)
% 
% \item Brücke (Pons)
% 
% \item verlängertes Mark (Medulla oblongata)
% 
% \item Ventrikelsystem
% 
% \end{itemize}
% \begin{itemize}
% \item Liquorgefüllte Hohlräume (Ventrikel)
% 
% \item haben nichts mit den Vnetrikel im Herz zu tun!
% 
% \end{itemize}
% Liquor = Gehirnwasser, in Ventrikeln gebildet

Das Hirn und das Rückenmark schwimmt in einer
Flüssigkeit, dem \gT{Liquor}. Die
Schädelknochen und der Liquor bieten eine Schutzfunktion. 
Wenn Liquor austritt (z.B. durch Nase, Ohr, ...)  ist das
ein Zeichen für eine schwerwiegende Verletzung!

\gFazit{Tritt Liquor in die Umwelt aus ist das ein Zeichen
für eine schwerwiegende Verletzung.} 

\gAuthNotes{KOCH: doppelt? GABS: ja. Schlüsselsätze sollen
doppelt (normales Format + Fazit-Box) vorkommen}

% \paragraph{\gZ{Blut-Hirn-Schranke}} 
% 
% \gZ{Die Blutgefäßwand ist im Hirn besonders dicht. Sie
% filtert daher besonders gut und nur für bestimmte Moleküle (\ce{O2},
% \ce{CO2}, Wasser, \ldots) leicht passierbar. Sie bietet dem Hirn
% Schutz vor schädlichen und giftigen Substanzen. 
% 
% Die Blut-Hirn-Schranke bietet dem Hirn besonderen
% Schutz vor schädlichen Substanzen im Blut.} 

\subsection{Das Periphere Nervensystem -- PNS}
\index{Nervensystem!Peripheres}
\index{PNS}\gTagNo{Unterteilung nach Lage und Aufbau}

\gDefinition{Unter dem Begriff \gT{Peripheres Nervensystem} fasst man
die Verbindungsnerven zwischen dem ZNS und den Endorganen zusammen.}

\gMarried
{\begin{itemize}
    \item Kommunikation ZNS ${\longleftrightarrow}$
    periphere  Organe
    \item Hirnnerven %(12)
    \item Spinalnerven (gehen vom Rückenmark segmental ab)
    \item elektrische Impulse, chemische Signale (Synapsen)
\end{itemize}

\paragraph{Arten von Nerven}~

\begin{itemize}
    \item motorisch \index{Nerven!motorische}  
    \item sensorisch \index{Nerven!sensorische} 
    \item vegetativ \index{Nerven!vegetative} 
    \begin{itemize}
        \item \gT{sympathisch}
        \item \gT{parasympathisch}
    \end{itemize}
\end{itemize}}
{\includegraphics[width=\linewidth]{./images/noauth/AASdev20090228-img34.png}
\gCaptionof{figure}{\gAuthNotesPicLegalNotOk{img34}{}Rückenmark
mit abgehenden peripheren Nerven. Beachte dass die Nerven
geordnet nach Rückenmarkssegmenten abgehen.
\gSourcePicture{Quelle nicht eruierbar.}}

}


\subsection{Das vegetative (autonome)  Nervensystem}
 \index{Nervensystem!Vegetatives} 
 \index{Nervensystem!autonomes}
 \index{Vegetatives Nervensystem}

 
\paragraph{Sympathikus  und Parasymphatikus  sind 
Gegenspieler}~
 \index{Sympathikus}  \index{Parasymphatikus} 
 \gTagNo{Unterteilung nach Funktion}

\gParallel{
\gT{Sympathikus}

Aktiviert den Aufmerksamkeits- und Kampfmodus.

\emph{"`Fight or flight"'} -- \emph{"`Kampf und Flucht"'}
}{
\gT{Parasympathikus} 

Aktiviert den Ruhemodus.

\emph{"`Rest and digest"'} - \emph{"`Rast und Verdauung"'}
}

z.B.: RR + HF

\gParallel{
Sympathikus:  HF ${\uparrow}$, RR ${\uparrow}$
}{
Parasymphatikus: HF ${\downarrow}$, RR ${\downarrow}$
}




\begin{figure}[H]
    \begin{center}
    \includegraphics[width=0.4\linewidth]{./images/noauth/AASdev20090228-img35.png}
    \includegraphics[width=0.4\linewidth]{./images/noauth/AASdev20090228-img36.png} 
    \gCaptionof{figure}{\gAuthNotesPicLegalNotOk{img35 
img36}{}Fight or flight. \gSourcePicture{Quelle nicht
eruierbar.}}
    \end{center}
\end{figure}

\begin{figure}[H]
    \begin{center}
   
    \includegraphics[width=0.8\linewidth]{./images/noauth/AASdev20090228-img37.png} 
    \gCaptionof{figure}{\gAuthNotesPicLegalNotOk{img37}{}
    Nur zur Illustration: Das vegetative Nervensystem hat
    Einfluss auf viele Organe. Dieses Bild
    \underline{nicht lernen}. Es dient nur der
    Verdeutlichung wie wichtig und einflussreich das vegetative
    Nervensystem im Körper ist. \gSourcePicture{Quelle nicht eruierbar.}} 
    \end{center}
\end{figure}
 
 
 
 
 
 
 
\clearpage
\section{Die Haut}
\index{Haut}
\label{topic:haut} 


Größtes Sinnesorgan und Grenzfläche zur
Außenwelt. \gZ{Sie macht ca. \nicefrac{1}{6} des
Körpergewichts aus und ihre Oberfläche beträgt beim
Erwachsenen ca. 1,6 m{\texttwosuperior}.}


\paragraph{Aufgaben}~

\begin{itemize}
    \item Schutz vor äußeren Einflüssen (mechanisch, Erreger,
    UV-Strahlung etc.)
    \item Schutz vor Wasserverlust/Austrocknung
    \item Temperaturregulation 
    \begin{itemize}
        \item Isolation: Schutz vor Wärmeverlust 
        \item Wärmeabgabe: Hautdurchblutung, Schweiß
\end{itemize}
    \item Sekretionsorgan (Schweiß, Talg)
    \item Sinnesorgan (Tastsinn, Temperatursinn)
\end{itemize}

\paragraph{Die Haut ist in Schichten aufgebaut}~
 
\gAuthNotes{Haut: Bild Querschnitt: muss vereinfacht werden,
lat. Ausdrücke weg bzw. grau}

\index{Oberhaut}\index{Lederhaut}\index{Unterhaut}

\gMarried{
\begin{itemize}
    \item \gT{Oberhaut} \gZ{(Epidermis)} \index{Oberhaut} 
    \begin{itemize}
        \item Hornschicht, geschichtetes Plattenepithel
    \end{itemize}
    \item \gT{Lederhaut} \gZ{(Corium, Dermis)}
    \index{Lederhaut}
    \begin{itemize}
        \item Hautanhangsgebilde
        \begin{itemize}
            \item Schweißdrüsen
            \item Talgdrüsen
            \item Haarfollikel, Nagelmatrix
        \end{itemize}
        \item Tastkörperchen, freie Nervenendigungen (Schmerz)
        \item \gZ{Kollagenfasern (Stütznetzwerk)}
        \item Arteriolen, Venolen (Versorgung der Haut, Temperaturregulation)
    \end{itemize}
    \item \gT{Unterhaut} (\gT{Subcutis}) \index{Unterhaut} 
    \index{Subcutis}
    \begin{itemize}
        \item Depotfett, Wärmeisolierung
    \end{itemize}
\end{itemize}
}{
\hspace{-0.5cm}\includegraphics[width=1.2\linewidth]{./images/noauth/AASdev20090228-img38.png}
\gCaptionof{figure}{\gAuthNotesPicLegalNotOk{img38}{Haut
Querschnitt}Die Haut im Querschnitt.
\index{Haut!Querschnitt}
\gSourcePicture{Quelle nicht eruierbar.}} }



\paragraph{Die Haut ist wichtig für die
\gEs{Thermoregulation}}~

Neben der Lunge dient vor allen Dingen  die Haut als wichtigstes Organ
zur \gEs{Thermoregulation}: Die Blutgefäße  der  Haut
sind stark ineinander verwickelt und können sich bei Bedarf stark erweitern. Die
stark durchblutete  Haut strahlt somit Wärme ab. Weiters wird durch
die Verdunstungskälte  des Schweißes  ebenfalls Wärme
abgegeben. 

\paragraph{Wundliegen}
\index{Wundliegen}\index{Dekubitus}\label{topic:dekubitus}Wenn
der Mensch nicht fähig ist in ausreichendem Umfang regelmäßig seine Lage zu ändern kann es zum wundliegen kommen. Die Wunden nennt man dann
\gT{Dekubitus}. \gZ{Besonders gefährdet sind jene Stellen,
an denen ein Knochen knapp unter der Haut liegt
(Druckstelle). Betroffen sind vor allem alte oder
pflegebedürftige bettlägrige Menschen, bewusstlose
Patienten bzw. sedierte Intensivpatienten sowie Patienten
mit Querschnittsläsionen.}


\clearpage
\section{Die Atemwege}
 \index{Atemwege} 

Die Atemwege gliedern sich in: 
\gAuthNotes{KOCH: mag Nasenrachenraum und Mundrachenraum
streichen, evt. wie Bild?}

\gMarried
{
\gT{Obere Atemwege}:

\begin{itemize}
    \item \gEs{Nasen}- und \gEs{Nasenhöhle} \index{Mundhöhle} 
    \index{Mundhöhle} 
    \item \gEs{Rachen} (gemeinsamer Weg von Nahrung und
    Atem!) \index{Rachen} 
    \item \gEs{Kehlkopf} (\gT{Larynx}) mit
    \gEs{Kehldeckel} \index{Kehlkopf} \index{Larynx} 
    \index{Kehldeckel}
\end{itemize}

}
{
\gT{Untere Atemwege}:
\begin{itemize}
    \item \index{Luftröhre}\gEs{Luftröhre}
    (\index{Trachea}\gT{Trachea})
    \item \index{Hauptbronchus}\gT{Hauptbronchus} (li./re.)
    \item \index{Bronchioli}\gT{Bronchioli}
    \item \index{Lungenbläschen}\gEs{Lungenbläschen}
    (\index{Alveolen}\gT{Alveolen})
\end{itemize}
}
%{\gTempPicInWork{Atemwege Übersicht img39}}

\begin{figure}[h]
    \includegraphics[width=\linewidth]{./images/anatomie/Respiratory_system_complete_de-edited.png} 
    
    \gCaptionof{figure}{\gAuthNotesPicLegalOk{}{PD}Die Atemwege, Übersicht.
\gSourcePicture{Mariana Ruiz Villarreal, Public domain}}
\label{fig:respiraktionstrakt-uebersicht}
\end{figure}



%\begin{multicols}{2}
\paragraph{Obere Atemwege}

Der Weg der Einatemluft beginnt in der Nase bzw. im Mund. Die Nase ist
mit Schleimhaut ausgekleidet, knöcherne Vorsprünge sorgen für
eine Vorwärmung der Luft. In der Schleimhaut be\-finden sich die
Sinneszellen für den Geruchssinn. 

\gZ{In den Schädelknochen befinden sich luft\-ge\-füllte
Hohlräume, die so\-genannten \gT{Neben\-höhlen} (Stirn-,
Kiefer\-höhlen, Siebbeinzellen). Diese stehen in Verbindung mit dem Nasenrachenraum.
Entzündet sich die Schleimhaut spricht man von der
\gE{Neben\-höhlen\-ent\-zündung}, welche oft langwierig
und schmerzhaft sein kann.} 

Die Luft gelangt nun in den \gT{Mund-Rachen\-raum} und den
\gT{Rachen}. \gEs{Hier kreuzt der Weg der Nahrung mit dem
Weg der Atemluft!} Die Nahrung gelangt vom Mundraum in die
hinten gelegene Speiseröhre, die Luft möchte vom Nasen\-rachen\-raum in die vorne gelegene Luftröhre. 

Um ein Verschlucken von Nahrung in die Lunge zu verhindern
verschließt der am Kehlkopf gelegene \gT{Kehldeckel}
während des Schluckaktes die Luftröhre. Dieser Reflex ist lebenswichtig! 

Im \gE{Kehlkopf}  befinden sich die \gT{Stimmbänder},
welche für die Lautbildung beim Sprechen wichtig sind.

\paragraph{Untere Atemwege}

In der \gT{Luftröhre} (\gT{Trachea}) wird die vom Rachen
und Kehlkopf kommende Luft zur Lunge weitergeleitet. Die Luftröhre teilt sich in einen
\gE{rechten} und einen \gE{linken} \gT{Hauptbronchus},
welche zu den jeweiligen Lungenflügeln ziehen. Jeder Hauptbronchus verzweigt sich immer weiter
bis schlussendlich sehr dünne \gT{Bronchioli} in die
\gT{Lungenbläschen} (\gT{Alveolen}) münden. Hier findet
der Gasaustausch statt.

Die Abschnitte der Atemwege, in denen kein Gasaustausch
stattfindet und der nur dem Transport dient, bezeichnet man
als {\quotedblbase}\index{Totraum}\gT{Totraum}{\textquotedblleft}.

%\end{multicols}


\subsection{Die Lunge ist für einen Teil des
Gasaustausches zuständig}
 \index{Gasaustausch!Lunge}  \index{Lunge}
 \index{Lunge!-nflügel}

Die Lunge hat \gEs{zwei Flügel}, welche sich jeweils in \gEs{3 (rechts)} bzw.
\gEs{2 (links) Lappen} gliedern: 

\begin{center}
\begin{tabularx}{\linewidth}[H]{XX}
\toprule
\multicolumn{2}{c}{\gTabHeading{Die Lunge}}
\\\toprule
\multicolumn{1}{c}{\gTabSub{Rechter Flügel}}
 &
\multicolumn{1}{c}{\gTabSub{Linker Flügel}}
\\
\begin{itemize}
    \item \gEs{3} Lappen
    \begin{itemize}
        \item Oberlappen \index{Oberlappen} 
        \item Mittellappen \index{Oberlappen} 
        \item Unterlappen \index{Unterlappen} 
    \end{itemize}
\end{itemize}
&
\begin{itemize}
    \item \gEs{2} Lappen
    \begin{itemize}
        \item Oberlappen
        \item Unterlappen
    \end{itemize}
\end{itemize}
\\\bottomrule
\end{tabularx}
\end{center}

Die Oberfläche ist mit dem \gT{Lungenfell}
\index{Lungenfell} \index{Pleura!lungenseitige} , der
\gE{lungenseitigen} \gT{Pleura}, überzogen. In einem Lungeflügel verzweigt sich der Hauptbronchus in immer weitere kleinere Bronchien bis schließlich die Bronchioli in
die \gT{Lungenbläschen} (\gT{Alveolen}) enden. 

\subsubsection{Die Alveolen sind die Endstücke
des Atemwegs und sind für den Gasaustausch zuständig}

\emph{Siehe Abb. \ref{fig:respiraktionstrakt-uebersicht}
rechts oben.}

In den \gT{Lungenbläschen} (\gT{Alveolen}) findet der
\index{Gasaustausch!Lunge}\gE{Gasaustausch} zwischen der
Einatemluft und dem Blut statt. Sauerstoff gelangt dabei aus den Alveolen in das Blut und
Kohlendioxid wird im Gegenzug vom Blut an die Alveolen
abgegeben.

Dabei besteht zwischen Luft und Blut kein direkter Kontakt, die Gase
müssen Hindernisse wie die Alveolenwand und die Blutgefäßwand
überwinden. \gE{Vergrößert sich dieser Weg aufgrund
einer Erkrankung ist die Sauerstoffaufnahme vermindert!} 

\gFazit{
Sauerstoff (\ce{O2}):   Alveolen ${\rightarrow}$ Blut

Kohlendioxid (\ce{CO2}):  Blut ${\rightarrow}$ Alveolen
}


\subsubsection{Die Pleura}\label{topic:pleura}

\gMarried{
\begin{itemize}
    \item \gE{Thoraxseitige} Pleura (\gT{Rippenfell}, fest mit
    Brustwand und Zwerchfell verwachsen, kleidet damit die
    Brusthöhle aus) \index{Rippenfell} 
    \index{Pleura!thoraxseitige}
    \item \gE{Lungenseitige} Pleura (\gT{Lungenfell},
    überzieht die Lungenoberfläche) \index{Lungenfell} \index{Pleura!lungenseitige}
    \item Aufgrund eines Unterdrucks haften die beiden Pleuren aneinander,
    sind aber gegeneinander verschieblich (wie 2 Glasplatten)
    \item Geringe Menge Pleuraflüssigkeit ${\rightarrow}$ Adhäsion der
    beiden Blätter!
    \item Normalerweise keine Verwachsungen der beiden Blätter
\end{itemize}
}
{\includegraphics[width=\linewidth]{./images/noauth/AASdev20090228-img41.png}
\gCaptionof{figure}{\gAuthNotesPicLegalNotOk{img41}{Pleura}Pleura.
\gSourcePicture{Quelle nicht eruierbar.}}}



\paragraph{Welche Funktion hat aber nun die Pleura?}

Da die eine Seite mit dem Brustkorb, die andere Seite mit der Lunge
verwachsen ist und der Raum zwischen den beiden Seiten luftdicht
abgeschlossen ist haften die beiden Seiten aneinander, aber sind
gegeneinander verschieblich. (Es gibt ja keine
Verwachsungen der beiden Seiten und die Pleuraflüssigkeit fungiert sogar ein
bisschen als "`Schmiermittel"'.) 

Wenn sich nun der Brustkorb oder das Zwerchfell bei der
Atmung bewegt, bewegt sich natürlich das Rippenfell mit (es ist ja mit dem Brustkorb
und dem Zwerchfell verwachsen). Da der Pleurazwischenraum
luftdicht ist, muss sich zwangsläufig aber auch das Lungenfell (an dem
wiederum die Lunge angewachsen ist) mitbewegen. Und wenn sich das Lungenfell
mitbewegt wird automatisch auch die Lunge wie eine Ziehharmonika
"`aufgezogen"' und Luft
"`angesaugt"'. 

\subparagraph{Frage:} \emph{Was würde passieren wenn der
Zwischenraum zwischen den Pleurablättern aufgrund einer Verletzung (z.B. ein
Loch aufgrund einer Stichwunde im Brustbereich) nicht mehr luftdicht
ist?} Antwort: s. Fußzeile (\footnote{Antwort: Die Lunge
würde in sich zusammenfallen, da sie ja nicht mehr aufgespannt wird. Siehe
"`Pneumothorax"'. \index{Pneumothorax!anatomische
Erklärung}})



\subsection{Atemmechanik}
\label{topic:atemmechanik} \index{Atemmechanik} 

Wenn man sich den Thorax annähernd als Zylinder vorstellt,
so kann eine Volumsvergrößerung (Einatmung) durch eine
Vergrößerung des Querschnittes (Heben der
Zwischenrippenmuskulatur\index{Zwischenrippenmuskulatur})
oder Vergrößerung der Höhe (Anspannung des Zwerchfells) erfolgen.

\begin{equation}
V = A \times h
\end{equation}
\begin{equation}
\mbox{Volumen} = \mbox{Fläche} \times \mbox{Höhe}
\end{equation}

\gAuthNotes{KOCH: entweder wird das genauer erklärt wie die
Funktion der Pleura oder es wird ganz einfach mündlich
erklärt.

GABS: Sollte schon schriftlich erklärt werden \ldots wir
wollen ja auch immer dass sie die "`Atemhilfsmuskulatur"'
etc. unterstützen, dann sollten sie wenigstens auch
nachlesen können "`wie man richtig atmet"'

KOCH: soll erklärt werden.}

\includegraphics[width=\linewidth]{./images/noauth/AASdev20090228-img42.png}
\gCaptionof{figure}{\gAuthNotesPicLegalNotOk{img42}{Atemmechanik}Atemmechanik.
\gSourcePicture{Quelle nicht eruierbar.}}

Für die Atemmachanik ist ein gutes Zusammenspiel von Skelett, Muskeln, Gewebe
und Pleura(spalt) wichtig. 

\gFazit{Ein bis zum Hals Verschütteter erstickt.}

\paragraph{Atemmuskulatur} Der wichtigste Atemmuskel ist
das \gEs{Zwerchfell}\index{Zwerchfell}. In Ruhe erledigt es
fast die gesamte Atemarbeit. Bei verstärkter Atmung kommen die Zwischenrippenmuskeln zum Einsatz \cite{LippertAnatomie7}.

\paragraph{Atemhilfsmuskulatur} \index{Atemhilfsmuskulatur}
Neben den Hauptatemmuskeln (Zwerchfell, Zwischenrippenmuskulatur) gibt es noch eine Reihe weiterer Atemhilfsmuskeln die
die Atmung unterstützen. Diese verbinden den Thorax
und die Oberarme. Fixiert man die Arme, z.B. durch
aufstützen, könne die Muskeln die Hebung des
Brustkorbes unterstützen. Vgl. dazu \emph{'Asthma', \ref{topic:asthma-bronchiale}} und \emph{'COPD', \ref{topic:copd}}.

\paragraph{Pleura} Die Funktion der Pleura ist unter \ref{topic:pleura} erklärt.

\gZ{Thoraxwand und Lunge besitzen elastische Eigenschaften: Sie werden gedehnt
und entwickeln elastische Rückstellkräfte. An der Lunge können zwei Mechanismen ein
Zusammenfallen des Lungengewebes bewirken. Erstens sind im Lungenzwischengewebe
reichlich elastische Fasern eingelagert, die die Eigenelastizität der Lunge
bewirken. Durch die Koppelung von Lunge und
Thoraxwand am Pleuraspalt entsteht ein elastisches Gesamtsystem. 

Bei entspannter
Atemmuskulatur und offenem Kontakt mit der Umgebungsluft wird sich ein
Gleichgewicht der elastischen Kräfte einstellen; der Brustkorb befindet sich in
der Atemruhelage. Durch den Zug der Lunge, die sich verkleinern möchte, und den
Gegenzug der Thoraxwand entsteht ein ständiger Unterdruck im Pleuraspalt.}

\subparagraph{Atemzyklus} 

\gZ{Volumenschwankungen des Brustraums bewirken
gleichzeitig auch Volumenschwankungen im Inneren der Lunge. Eine Vergrößerung des Thoraxinnenraums
wird durch Kontraktion der Atemmuskulatur bewirkt. Dieser Vorgang wird als aktive
Phase der Atmung bezeichnet. Dabei kommt es vor allem zu einer Abflachung und
einem Tiefertreten des Zwerchfells, aber auch zu einer Anhebung der Rippen durch
die äußere Zwischenrippenmuskeln. Als Folge entsteht in der Lunge im Vergleich
zur Umgebungsluft ein Unterdruck, der durch Einströmen von Luft ausgeglichen
wird. Thoraxwand und Lunge werden dabei gedehnt. 

Am Ende einer ruhigen Einatmung
erschlafft die Einatmungsmuskulatur. Die elastischen Rückstellkräfte führen das
Thoraxvolumen wieder in die Ausgangslage zurück. Dabei kommt es zu einer
Druckerhöhung in der Lunge, die eine Luftbewegung hin zur Umgebungsluft bewirkt.
Diese erfolgt in aller Regel passiv, d. h. ohne aktive
muskuläre Anstrengung. Dieser Vorgang erfolgt in Ruhe bei einem Erwachsenen etwa zwölfmal in der Minute
und wird als Atemfrequenz bezeichnet.}


\gFazit{Der wichtigste Atemmuskel ist das \gEs{Zwerchfell}!}

\subsection{\gZ{Atemfrequenz, -tiefe und -volumen}}

Die Atemfrequenz, die Atemtiefe und das Atemvolumen
werden im Kapitel Vitalfunktionen unter
\gSee{topic:atemvolumina} besprochen.

\subsection{Luft}

Die Luft, welche wir atmen, ist ein Gasgemisch aus
Stickstoff, Sauerstoff, Kohlendioxid und Edelgasen.
\index{Stickstoff}  \index{Sauerstoff!Raumluftgemisch} 
\index{Kohlendioxid!Raumluftgemisch} Zwischen der normalen
Einatem- und Ausatemluft besteht besonders hinsichtlich des Sauerstoff- und Kohlendioxid-Anteils ein Unterschied, siehe
Tab. \ref{tab:luft-gas}.

\begin{table}[H]
\gCaptionof{table}{Gasgehalt in der Raum- und
Ausatemluft.\label{tab:luft-gas}}

\begin{tabular}[H]{ccc}
\addlinespace
\centering \gTabHeading{Raumluft} & & \gTabHeading{Ausatemluft}
\\\addlinespace\midrule\addlinespace
\gEs{21\,\%}  & Sauerstoff (\ce{O2}) & \gEs{16\,\%}
\\\addlinespace
78\,\% & Stickstoff (\ce{N}) & 78\,\%
\\\addlinespace
\gEs{0,03\,\%} & Kohlendioxid (\ce{CO2}) & \gEs{4\,\%} 
\\\addlinespace
1\,\% & Edelgase  & 1\,\%
\\\addlinespace\bottomrule
\end{tabular}
\end{table}


\clearpage
\section{Der Verdaungstrakt -- Von der Schüssel in die
Schüssel}
\index{Magen-Darm-Trakt}
\index{Gastrointestinaltrakt}
\index{GIT}
\label{topic:gastrointestinaltrakt}

\emph{Syn.: \gT{Magen-Darm-Trakt},
\gT{Gastrointestinaltrakt} (\gT{GIT})}


\begin{itemize}
    \item Die Verdauung dient dem Stoffwechsel.
    \item Prozess:
    \begin{itemize}
        \item Nahrungsaufnahme
        \item Weiterverarbeitung im Magen-Darm-Trakt
        \item Transport der Nährstoffe und Energieträger  ins Blut
        \item Zellaufbau und Abbau von Abfallstoffen
    \end{itemize}
    \item Über die Verdauung gewinnt der Körper  also einerseits die
    nötigen {\quotedblbase}Bausteine{\textquotedblleft},  andererseits
    auch die nötige Energie um  Zellen aufbauen und in Funktion erhalten
    zu können.
\end{itemize}

\paragraph{Übersicht}~

\begin{figure}[H]
    \includegraphics[width=0.8\linewidth]{./images/anatomie/Digestive_system_diagram_de.png}
    \gCaptionof{figure}{\gAuthNotesPicLegalOk{}{}Schema des Verdauungstraktes.
    \gSourcePicture{Mariana Ruiz Villarreal, Public domain}}
\end{figure}


\begin{gWideParEnv}
\begin{minipage}{\linewidth}
\gCaptionof{table}{Übersicht über den Verdauungstrakt --
\emph{"`Von der Schüssel in die Schüssel."'}\label{tab:verdauungstrakt-uebersicht}}
\begin{tabulary}{\linewidth}[H]{LLL}
\toprule\addlinespace
\noindent\multirow{3}{8em}{\gTabHeading{Mundhöhle}}
 &
 Zunge
 &
\noindent\multirow{2}{15em}{Zerkleinerung der Nahrung und Vermischung
mit Speichel}
\\\addlinespace\hhline{~-~}\addlinespace
 &
Zähne 
&
\\\addlinespace\hhline{~-~}\addlinespace
 &
Speicheldrüsen
&
\\\addlinespace\hhline{---}\addlinespace
Rachen
 &
 &
Kreuzung des Speisebreis mit dem Luftweg

Kehldeckel verschließt beim Schlucken die Luftröhre
\\\addlinespace\hhline{---}\addlinespace
\noindent\gEs{Speiseröhre} (\gZ{Ösophagus})
 &
 &
Speisebrei wird mittels Muskelbewegung zum Magen befördert
\\\addlinespace\hhline{---}\addlinespace
\noindent\gTabHeading{Magen}
 &
 &
Vermischung mit Magensäure

Magensäure zerlegt Nahrung

Pförtner gibt Portionsweise Speisebrei an den Zwölffingerdarm ab

\\\addlinespace\hhline{---}\addlinespace
\multirow{4}{8em}{\gTabHeading{Dünndarm}}
&
\noindent\gTabSub{Zwölffingerdarm} \gZ{Duodenum}
 &
Gallenflüssigkeit (${\leftarrow}$Leber) und Bauchspeichel
(${\leftarrow}$Pankreas) wird zugesetzt
\\\addlinespace\hhline{~--}\addlinespace
 &
\noindent\gTabSub{Krummdarm} \gZ{Jejunum}
 &
Spaltung von Nährstoffen (Eiweiß, Kohlehydrate, Fett) und
Nährstoffaufnahme
\\\addlinespace\hhline{~-~}\addlinespace
 &
\noindent\gTabSub{Leerdarm} \gZ{Ileum}
 &
\\\addlinespace\hhline{---}\addlinespace
\multirow{6}{8em}{\gTabHeading{Dickdarm}\\(\gTabHeading{Colon})}
&
\noindent\gTabSub{Blinddarm}
&
Sackgasse
\\\addlinespace
&
\noindent\gTabSub{Wurmfortsatz} (\gTabSub{Appendix}) 
 &
Anhängsel der Sackgasse, kann sich entzünden.
\\\addlinespace\hhline{~--}\addlinespace
 &
\noindent\gTabSub{Aufsteigender} Teil
 &
\multirow{4}{\columnwidth}{Wasserentzug, \\Eindickung, \\
Bakterienbesiedlung, \\Produktion und Speicherung des
Stuhls} 
\\\addlinespace\hhline{~-~}\addlinespace
 &
\noindent\gTabSub{Querender} Teil
 &
\\\addlinespace\hhline{~-~}\addlinespace
 &
\noindent\gTabSub{Absteigender} Teil
 &
\\\addlinespace\hhline{~-~}\addlinespace
 &
\noindent\gTabSub{S-förmig} gebogener Teil (\gZ{Sigma})
 &
\\\addlinespace\midrule\addlinespace
\noindent\gTabHeading{Mastdarm} (\gTabHeading{Rektum})
 &
 &
{\quotedblbase}Speicherung{\textquotedblleft} des ausscheidungsbereiten
Stuhls
\\\addlinespace\hhline{---}\addlinespace
\noindent\gTabHeading{Anus}
 &
 &
Absetzen des Stuhles

\\\addlinespace\bottomrule
\end{tabulary}

\end{minipage}
\end{gWideParEnv}


\subsection{Mundhöhle}\index{Mundhöhle}

\begin{description}
    \item[Zähne] zerkleinern die Nahrung
    \item[Speicheldrüsen] produzieren Speichel um Nahrung
    gleitfähig zu machen und Nahrung mit Verdauungsenzymen \gZ{(für Zucker und Stärke)} zu versetzen.
    \item[Zunge] schiebt Nahrung in Schlund weiter,
    verschließt beim Schlucken Kehldeckel (sonst  \gEs{Aspiration!})
\end{description}



\subsection{Die Speiseröhre (Ösophagus)}
\index{Speiseröhre}\index{Oesophagus@Ösophagus} 

\gMarried
{
\begin{itemize}
    \item ca. 24cm langer, dreischichtiger Muskelschlauch
    \gZ{(oberes Drittel quergestreifte Muskulatur, untere
    \nicefrac{2}{3} glatt)}
    \item Verbindung Mund $\rightarrow$ Magen
    \item "`schiebt"' Nahrung durch
    \gEs{Muskelkontraktionen}
    (\gT{Peristaltik}\index{Peristaltik}) weiter
    \item hier findet keine Verdauung statt
\end{itemize}
}
{
\centering
\includegraphics[width=0.3\linewidth]{./images/noauth/AASdev20090228-img44.png} 
\includegraphics[width=0.3\linewidth]{./images/noauth/AASdev20090228-img45.png}
\gCaptionof{figure}{\gAuthNotesPicLegalNotOk{img44
img45}{}Ösophagus. \gSourcePicture{Quelle nicht eruierbar.}}
}




\subsection{Der Magen}
\index{Magen}

\gZ{(Gaster, Ventriculus)}

\gAuthNotes{KOCH: Fließtext!}
\begin{itemize}
    \item 2-3\,l Magensaft
    \item Schleim (schützt  Magenschleimhaut)
    \item Salzsäure (Bakterienfalle, pH
    1,0)\index{Salzsäure}
    \item \gZ{Pepsin  (Enzym für Eiweiß-Verdauung)}
    \item Wandmuskulatur \gEs{durchmischt} den Speisebrei
    mit der Magensäure
    \item portionsweise Abgabe von Speisebrei durch den
    Magenpförtner \index{Magen!-pförtner} \gZ{(Pylorus)} in
    den Zwölffingerdarm
\end{itemize}


\gParallel{
\centering\includegraphics[width=0.7\linewidth]
{./images/noauth/AASdev20090228-img46.png} 
\gCaptionof{figure}{\gAuthNotesPicLegalNotOk{img46
img45}{}Lage des Magens. \gSourcePicture{Quelle nicht
eruierbar.}}
}{
 \includegraphics[width=\linewidth]{./images/noauth/AASdev20090228-img47.png} 
 \gCaptionof{figure}{\gAuthNotesPicLegalNotOk{img47}{}Der
 Magen, der Pförtner und das erste Stück
 des Zwölffingerdarms im Durchschnitt.
 \gSourcePicture{Quelle nicht eruierbar.}}
}




\subsection{Zwölffingerdarm (Duodenum)}

\begin{itemize}
\item Beginn des Dünndarmes, ca 25\,cm lang
\item Verdauungsenzyme aus Leber und Pankreas für Fett, Eiweiß,
Zucker münden hier ein
\end{itemize}

\gParallel{
\includegraphics[width=\linewidth]{./images/noauth/AASdev20090228-img48.png}
\gCaptionof{figure}{\gAuthNotesPicLegalNotOk{}{}Die Lage des
Zwölffingerdarms im Körper.} }{
\includegraphics[width=\linewidth]{./images/noauth/AASdev20090228-img49.png}
\gCaptionof{figure}{\gAuthNotesPicLegalNotOk{}{}Der
Zwölffingerdarm und die Bauchspeicheldrüse.} }

\gAuthNotes{KOCH: sind beides wirklich Enzyme?

GAB: Ja.}


\subsection{Die \index{Leber}Leber} 

\gAuthNotes{KOCH: fehlt: was ist sie, eine Drüse, ein
Hohlmuskel, \ldots } 

\gZ{("`Hepar"')}

Die Leber liegt im rechten Oberbauch und hat einen linken und
    rechten Lappen.
    
\gMarried
{
\gSlideHeading{Aufgaben}

\begin{itemize}
\item "`Chemiefabrik"' des Körpers: Produktion von vom
    Körper benötigten Eiweißen
    \begin{itemize}
        \item Produktion von vom
    Körper benötigten Eiweißen
    \item Entgiftung: Abbau von Giftstoffen (die Leber nimmt
    durch die Giftstoffe mit der Zeit selbst auch Schaden
    (Leberentzündung, Leberzirrhose)
    \item erhält über \index{Pfortader}\gT{Pfortader}
    (Vena portae) nährstoffreiches Blut vom Darm zwecks Entgiftung   
\end{itemize}
    \item {\quotedblbase}Verdauungsdrüse{\textquotedblleft}
    \begin{itemize}
        \item produziert \index{Gallenflüssigkeit}Gallenflüssigkeit
        ({\quotedblbase}Galle{\textquotedblleft}), welche in der Gallenblase
        gespeichert wird: Fettverdauung
        \item Von der Leber gehen die Gallengänge, an den Gallengängen
        hängt die
    \end{itemize}
    \item Speicherung von Fett \& Zucker
    
\end{itemize}}
{
\includegraphics[width=\linewidth]{./images/noauth/AASdev20090228-img50.png}
}

\subsubsection{Gallenblase}
\index{Gallenblase}
\label{topic:gallenblase}

Am Gallengang angehängt findet sich die \gT{Gallenblase}.
Sie dient nur als Zwischenspeicher für die Gallenflüssigkeit.

\gAuthNotes{GABS: Erklärung/Bild Pfortader einfügen? fragt
STEU}




\gDias{Bilderserie: Leber und Gallenblase.}
{\includegraphics[width=\linewidth]{./images/noauth/AASdev20090228-img51.png}
\gCaptionof{figure}{\gAuthNotesPicLegalNotOk{img51}{}Leber
Ansicht von vorne.
\gSourcePicture{Quelle nicht eruierbar.}}}
{\includegraphics[width=\linewidth]{./images/noauth/AASdev20090228-img52.png}
\gCaptionof{figure}{\gAuthNotesPicLegalNotOk{img52}{}Leber,
Ansicht von unten mit Gallenblase.
\gSourcePicture{Quelle nicht eruierbar.}}}
{\includegraphics[width=\linewidth]{./images/noauth/AASdev20090228-img53.png}
\gCaptionof{figure}{\gAuthNotesPicLegalNotOk{img53}{}Gallenblase.
\gSourcePicture{Quelle nicht eruierbar.}}}


 

\subsection{Bauchspeicheldrüse}\index{Bauchspeicheldrüse}

(\index{Pankreas}\gT{Pankreas})

Das Pankreas ist ca.\,15-22\,cm lang und liegt hinter dem
Magen. Es hat einen \gEs{Ausführungsgang}, welcher
zusammen mit dem Gallengang \gEs{in den Zwölffingerdarm}
(Duodenum) mündet.

\paragraph{Aufgaben} 

Das Pankreas hat \gEs{zwei} wesentliche Funktionen:

\begin{description}
    \item[Enzymproduktion] für Verdauung
    ${\rightarrow}$ Zwölffingerdarm
    \item[Hormonproduktion] Ausschüttung ins Blut
    \begin{itemize}
        \item \index{Insulin}\gT{Insulin} ${\rightarrow}$ Blutzucker
        ${\downarrow}$
        \item \gZ{Glukagon ${\rightarrow}$ Blutzucker ${\uparrow}$}
    \end{itemize}
\end{description}

%\gTempPicInWork{Gallensystem}

\subsection{Krumm- und Leerdarm -- Jejunum und Ileum}
\index{Krummdarm}\index{Leerdarm}\index{Jejunum}\index{Ileum}

\gT{Jejunum}, \gT{Ileum}

\gMarried
{\begin{itemize}
    \item ca. 3\,m lang
    \item Oberflächenvergrößerung durch Darm\-zotten
    \item \gEs{Nährstoffaufnahme}\index{Nährstoffaufnahme}
    \item \gZ{beweglich, am Darmgekröse (Mesenterium)
    aufgehängt}
\end{itemize}}
{
\includegraphics[width=0.6\linewidth]{./images/noauth/AASdev20090228-img56.png}
\gCaptionof{figure}{\gAuthNotesPicLegalNotOk{img56}{}Dünndarm
Querschnitt.} }





\subsection{Dickdarm (Colon)}
\index{Dickdarm}\index{Colon}

Der Dickdarm (Colon) ist ca. 1,5\,m lang und wird in
verschiedene Abschnitte unterteilt: 

\begin{center}
\begin{tabulary}{\linewidth}[H]{llL}
\toprule
\multirow{6}{8em}{\gEs{Dickdarm}\\(\gT{Colon})}
&
\gEs{Blinddarm}\index{Blinddarm}
&
Sackgasse
\\
&
\gEs{Wurmfortsatz} (\gT{Appendix})
\index{Wurmfortsatz}\index{Appendix} &
Anhängsel der Sackgasse, kann sich entzünden.
\\\hhline{~--}
 &
\gEs{Aufsteigender} Teil \index{Dickdarm!aufsteigender
Teil} &
\multirow{4}{\columnwidth}{Wasserentzug, Eindickung, \\
Bakterienbesiedlung, \\Speicherung} 
\\\hhline{~-~}
 &
\gEs{Querender} Teil \index{Dickdarm!querender 
Teil}
 &
\\\hhline{~-~}
 &
\gEs{Absteigender} Teil \index{Dickdarm!absteigender
Teil}
 &
\\\hhline{~-~}
 &
\gEs{S-förmig} gebogener Teil (\gT{Sigma})
\index{Dickdarm!S-förmig gebogener Teil} &
\\\bottomrule
\end{tabulary}
\end{center}

\gMarried{
Der Leerdarm mündet in den Dickdarm. Diese Mündung ist wie eine
\gE{T-Kreuzung} gebaut: \gE{Nach oben} geht es in den
\gEs{aufsteigenden Teil} des Dickdarmes, \gE{nach unten} in
den \gEs{Blinddarm}. Wie der Name des Blinddarmes schon sagt
endet er blind, er ist daher eine Sackgasse. Am Blinddarm
angehängt findet man den \gEs{Wurmfortsatz}
(\gT{Appendix}) welcher sich häufig entzündet
(\gT{Appendiz\underline{itis}\footnote{Das Anhängsel
\emph{"`-itis"' zeigt immer eine Entzündung an, vgl.
Kapitel Hygiene}}}). Der Wurmfortsatz erfüllt beim
Menschen keine nennenswerte Funktion, außer der
Arbeitsplatzsicherung von Anästhesisten und Chirurgen
(wegen der vielen Blinddarm-Operationen).
}{\includegraphics[width=0.8\linewidth]{./images/noauth/AASdev20090228-img57.png}

\gCaptionof{figure}{\gAuthNotesPicLegalNotOk{img57}{}Lage
des Dickdarms und des Appendix. \gSourcePicture{Quelle
nicht eruierbar.}
}
}





\clearpage
\section{Sonstige Bauchorgane}
\label{topic:sonstige-bauchorgane}

\subsubsection{Bauchfell}
\index{Bauchfell}

\begin{itemize}
    \item Die gesamte Bauchhöhle ist mit  Bauchfell
    (\index{Peritoneum}Peritoneum) ausgekleidet
    \item dünner Überzug
    \item {\quotedblbase}Aufhängung{\textquotedblleft} der Bauchorgane
    \item zusätzlich großes und kleines Netz
    \item Kann sich Entzünden (\gT{Bauchfellentzündung}
    =
    \gT{Peritonitis}\index{Peritonitis}\index{Bauchfellentzündung}): zB bei Darmverletzung, Bauchfelldialyse
\end{itemize}


\gMarried
{Der Bauch "`aufgeklappt"': Nach oben geklappt sieht man das
 "`Große Netz"', auf die Seiten
und nach unten die Bauchmuskulatur. 

Die hellen, fleischfarbenen Schlingen sind Dünndarmschlingen, die
dicken braunen Schlingen sind der aufsteigende und der quere Teil des
Dickdarmes.}
{
\includegraphics[width=\linewidth]{./images/noauth/AASdev20090228-img58.png}
\gCaptionof{figure}{\gAuthNotesPicLegalNotOk{}{}Der Bauch, geöffnet, mit aufgeklapptem
Bauchfell.}}


\subsection{Die Milz} 
\index{Milz}

\gZ{Syn: Lien, Splen.}



\gMarried{
\begin{itemize}
\item hat nichts  mit der Verdauung  zu tun 
\item 12\,cm, oval, von \gEs{Kapsel}\index{Kapsel!Milz}
umgeben
\item im linken Oberbauch unter den Rippen gelegen
\item Teil des Immunsystems
\item entfernt alte Erythrozyten, beim Embryo auch Blut\-bildung
\item sehr gut durchblutet, Blutspeicher (Gefahr:
\index{Milzriss}Milzriss)
\end{itemize}
}{
\centering\includegraphics[width=0.6\linewidth]{./images/noauth/AASdev20090228-img60.png}
\gCaptionof{figure}{\gAuthNotesPicLegalNotOk{img60}{}Milz:
Lage.
\gSourcePicture{Quelle nicht eruierbar.}}
}



\gCave{\underline{Achtung}: sog. {\quotedblbase}\index{Zweizeitiger
Milzriss}\index{Milz!-riss, zweizeitiger}\gT{Zweizeitiger
Milzriss}{\textquotedblleft}:

Zuerst reißt die Milz ein, die Kapsel bleib erhalten, es blutet zuerst in die Kapsel und der Druck steigt.
Nach einer gewissen Zeit reißt auch die Kapsel und es blutet nun
ungehindert aus der inneren Wunde. \gEs{Der Patient wirkt
zuerst {\textquotedblleft}ganz normal{\textquotedblleft}, unerwartet und
plötzlich verfällt er.}} 



\subsection{Die Niere(n) (Ren)}
\index{Niere}

\gMarried{
\begin{itemize}
    \item paarig angelegtes Organ
    \item stark durchblutet, empfindlich auf RR- Abfall ${\rightarrow}$
    Funktion ${\downarrow}$
    \item reguliert
    \begin{itemize}
        \item \gEs{Wasser-} und
        \gEs{Elektrolythaushalt}
        \index{Wasserhaushalt!Niere}
        \index{Elektrolythaushalt!Niere}
        \item Volumshaushalt\index{Volumshaushalt!Niere}
        \item pH
        (Säure/Basen-Haushalt)
        \index{Säure-Basen-Haushalt!Niere}
    \end{itemize}
    \item Geht im Alter, besonders bei Diabetikern, kaputt ${\rightarrow}$
    {\quotedblbase}\gT{Chronische
    Niereninsuffizienz}{\textquotedblleft} ${\rightarrow}$
    {\quotedblbase}\gT{Blutwäsche}{\textquotedblleft}
    (\index{Dialyse}\gT{Dialyse})
\end{itemize}
}{
\centering\includegraphics[width=0.8\linewidth]{./images/noauth/AASdev20090228-img61.png}
\gCaptionof{figure}{\gAuthNotesPicLegalNotOk{img61}{}Niere.
\gSourcePicture{Quelle nicht eruierbar.}}
}



\paragraph{Aufbau}~

\begin{itemize}
\item Rinde, Mark
\item Hilus (Stiel)
\item Gefäße
\item Nierenbecken mit Anschluss an den
Harntrakt\index{Nierenbecken}
\end{itemize}

\gAuthNotes{Niere: Anatomie und Funktion muss noch
ausgearbeitet werden.}


\gAuthNotes{KOCH: grau:}
\paragraph{\gZ{Mikroskopischer Aufbau}}

\gZ{
\begin{itemize}
    \item Nephron ist die funktionelle Einheit (ca. 1
    Mio/Niere)
    \begin{itemize}
        \item Glomerulus: Hier wird das Blut gefiltert und Plasmawasser
        abgepresst (ca. 150\,l\,/\,Tag)
        \item Schleifenapparat (Tubulus): Das abgepresste Plasma\-wasser wandert
        den Schleifen\-apparat entlang, hier werden die meisten wichtigen
        Stoffe und der Groß\-teil des Wassers wiederum zurückgeholt,
        sodass am Ende nur mehr 1,5 -- 2~l Harn pro Tag in der Harnblasen
        landen.
    \end{itemize}
\end{itemize}
}

\gZ{
\paragraph{Funktionsweise}

\begin{itemize}
    \item 1500 l Blut pro Tag wird
    {\quotedblbase}bearbeitet{\textquotedblleft}
    \item Glomerula: Filtration ${\rightarrow}$ 150\,l
    \index{Primärharn}Primärharn (enthält Elektrolyte, etc.) wird
    gewonnen und fließt durch den Schleifenapparat
    \item Schleifenapparat: {\quotedblbase}Recycling{\textquotedblleft},
    Gutes kommt zurück ins Blut, der Rest fließt weiter Richtung
    Harnblase ${\rightarrow}$ 1,5-2\,l (End-)Harn/Tag
\end{itemize}
}

\subsection{Die Nebenniere(n) (Adren)}

% \gTempPicMissing{Nebennieren}

\gZ{Paariges, jeweils am oberen Nierenpol gelegenes Organ.

Hat nichts  mit dem Harntrakt oder der eigentlichen Nierenfunktion zu
tun.

\gEs{Hormonproduktion}
\gZ{
\begin{itemize}
    \item Cortison
    \item Sexualhormone (Androgene)
    \item Adrenalin , Noradrenalin
\end{itemize}}

Die Nebennieren sitzen wie {\quotedblbase}Mützen{\textquotedblleft}
auf den Nieren, haben mit deren Funktion allerdings
\gE{nichts} zu tun.
}

\subsection{Harnableitendes System -- Harntrakt}
\index{Harn!-ableitendes System}\index{Harn!-trakt}

\gMarried
{Der Harntrakt besteht aus 
\begin{itemize}
    \item Nieren inkl. Nierenbecken
    \item \index{Harn!-leiter}je 1 Harnleiter 
    \gZ{(Ureter)}, 30\,cm lang, mündet  in die 
    \item \index{Blase}Blase \index{Harn!-blase}
    \item Harnröhre  (\gZ{Urethra,} Mann: 25\,cm, Frau
    4\,cm)\index{Harn!-röhre}
\end{itemize}

Die Harnröhre verläuft beim Mann von der Harnblase durch die
Prostata \index{Harn!-röhre!Verlauf} und den Penis und
mündet an der Eichel in die Außenwelt. Bei einer Vergrößerung der Prostata kann der Harnabfluss bis
zur Unpassierbarkeit erschwert sein.  

Bei der Frau ist die Harnröhre kürzer als beim Mann und
mündet oberhalb des Scheideneinganges. 

\gCaptionof{figure}{\gAuthNotesPicLegalOk{}{Lena}Harntrakt
mit Nieren (rechte Niere im Querschnitt mit Nierenbecken),
Harnleiter, Blase und Harnröhre. \label{fig:harntrakt}
\gSourcePicture{Lena Hirtler}}
}
{~

\includegraphics[width=\linewidth]{./images/anatomie/anatomie-harntrakt-edited.png}
} 





\subsection{Was gibt es noch?}

Außer den "`klassischen Bauchorganen"' liegen auch andere
Organsysteme im Bauchraum auf die nicht vergessen werden
sollte. Besonders betrifft das die weiblichen
Fortpflanzungsorgane wie Eierstöcke, Eileiter und Uterus.
Diese werden allerdings im Kapitel "`Geschlechtsorgane"'
gesondert besprochen. 

\clearpage
\section{Geschlechtsorgane}
\subsection{Männliche Geschlechtsorgane}

\begin{itemize}
    \item innen:
    \begin{itemize}
        \item Hoden \index{Hoden}
        \begin{itemize}
            \item Spermienproduktion
            \item Geschlechtshormon-Produktion
        \end{itemize}
        \item Nebenhoden: Spermienreifung und Lagerung
        \item Samenstrang, besteht aus
        \begin{itemize}
            \item Samenleiter \index{Samenleiter}
            \item Nerven
            \item Gefäße
            \item durchzieht Leistenkanal
        \end{itemize}
        \item \index{Prostata}\gT{Prostata}
        (Vorsteherdrüse)
        \begin{itemize}
            \item produziert Samenflüssigkeit (Beweglichkeit der Spermien)
            \item \gEs{Vergrößert sich im Alter} und
            kann die \gEs{Harnröhre abdrücken}, es kommt
            zum Harnstau ${\rightarrow}$ häufiger Notfall!
        \end{itemize}
    \end{itemize}
    \item außen:
    \begin{itemize}
        \item Glied (Penis): Schwellkörper, Harnsamenröhre, Begattungsorgan
        \item Hodensack \gZ{(Skrotum)}: umhüllt Hoden und
        Nebenhoden
        \begin{itemize}
            \item Hodenstiel kann sich verdrehen
            \index{Hodenstiel} ${\rightarrow}$ \gCa{Notfall!}
        \end{itemize}
    \end{itemize}
\end{itemize}

\begin{figure}[H]
    \begin{center}
     \includegraphics[width=7cm]{./images/anatomie/anatomie-genitale-maennlich-edited.png}
    \gCaptionof{figure}{\gAuthNotesPicLegalOk{}{}Die
    männlichen Geschlechtsorgane im
    Querschnitt. \gSourcePicture{Lena Hirtler}}
    \end{center}
\end{figure}
 

\subsection{Weibliche Geschlechtsorgane}

\begin{itemize}
    \item innen:
    \begin{itemize}
        \item \gEs{Eierstöcke} \gZ{(Ovar)}: Reifung der
        Eizellen \index{Eierstöcke}
        \item \gEs{Eileiter} \gZ{(Tuben)}: Transport des
        reifen Eis in Richtung Gebärmutter \index{Eileiter}
        \item \gEs{Gebärmutter} \index{Gebärmutter}
        (\index{Uterus}\gT{Uterus}): dehnbarer Hohlmuskel,
        mit Schleimhaut ausgekleidet,
        \begin{itemize}
            \item Schleimhaut wird regelmäßig aufgebaut und wieder
            abgestoßen (Menstruation, Regelblutung)
            \item Hier nistet sich die befruchtete Eizelle
            ein.
            \item Der Gebärmutterhals
            \index{Gebärmutterhals}
            (\gT{Cervix}\index{Cervix}) führt in Richtung der Scheide.
            \item Die Mündung der Gebärmutter in die
            Scheide ist der \gT{Muttermund}.
            \index{Muttermund!Anatomie}
        \end{itemize}
        \item \gEs{Scheide}
        (\gT{Vagina})\index{Scheide}\index{Vagina}
    \end{itemize}
    \item außen:
    \begin{itemize}
        \item Schamlippen \gZ{(Labien)}:
        goße/äußere, kleine/innere
        \item \gT{Klitoris}\index{Klitoris}: entspricht der
        männlichen Eichel
        \item Öffnung der Harnröhre
    \end{itemize}
\end{itemize}


\begin{figure}[H]
    \begin{center}
    \includegraphics[width=7cm]{./images/anatomie/anatomie-genitale-weiblich-edited.png}
    \gCaptionof{figure}{\gAuthNotesPicLegalOk{}{Lena
    }Weibliche Geschlechtsorgane, Querschnitt.
    \gSourcePicture{Lena Hirtler}}
    \end{center}
\end{figure}

 

\subsubsection{Der weibliche Zyklus}
\label{topic:weiblicher-zyklus}\index{Zyklus, weiblicher}

\begin{itemize}
    \item Der weibliche Zyklus dauert 28 Tage
    \item gerechnet von Menstruation zu Menstruation  (Monatsblutung )
    \item Uterusschleimhaut wird aufgebaut und wieder abgestoßen
    \item hormongesteuert
    \begin{itemize}
        \item Östrogen, Progesteron
    \end{itemize}
    \item ca. am 14. Tag: Eisprung
    \item Befruchtung nein:
    \begin{itemize}
        \item Zyklus läuft normal weiter, Menstruation
    \end{itemize}
    \item Befruchtung ja:
    \begin{itemize}
        \item Eizelle nistet sich im Uterus ein, Zyklus gestoppt
    \end{itemize}
\end{itemize}

Was kann dabei schief gehen? Wenn sich die Eizelle schon im Ei\-leiter
einnistet kommt es zu einer
\index{Eileiter\-schwanger\-schaft}\gEs{Eileiter\-schwanger\-schaft}
die lebens\-bedrohlich ist (Eileiter reißt und die Frau verblutet).

\begin{minipage}{\linewidth}
\includegraphics[width=\linewidth]{./images/noauth/AASdev20090228-img65.png}
\gCaptionof{figure}{\gAuthNotesPicLegalNotOk{img65}{}Der weibliche Zyklus. Oben:
Follikelreifung. Unten: Auf- und Abbau der Uterusschleimhaut. Am Tag 0 des Zyklus wird die Uterusschleimheut abgestoßen, es kommt zur
Regelblutung, danach wird die Schleimhaut wieder aufgebaut. Am Tag 14
erfolgt der Eisprung. Kommt es zu keiner Befruchtung wird die
Schleimheut erneut am Tag 28 abgestoßen. Nach einer
Befruchtung wandert die Eizelle weiter in den Uterus und nistet sich dort ein.
\gSourcePicture{Quelle nicht eruierbar.}}
\end{minipage}



